\documentclass[11pt,a4paper]{report}

\usepackage[utf8]{inputenc}
\usepackage[T1]{fontenc}
\usepackage{textcomp}
\usepackage[french]{babel}
\addto\captionsfrench{\renewcommand{\tablename}{Tableau}
  \renewcommand{\figurename}{Figure}}

\usepackage{lmodern}\usepackage{geometry}\usepackage{setspace}
\usepackage{microtype}\usepackage{graphicx}\usepackage{xcolor}\usepackage{colortbl}
\usepackage{booktabs}\usepackage{tabularx}\usepackage{longtable}
\usepackage{array}\usepackage{multirow}\usepackage{enumitem}
\usepackage{titlesec}\usepackage{fancyhdr}\usepackage{caption}
\usepackage{placeins}\usepackage{float}\usepackage{needspace}
\usepackage{amssymb}\usepackage{amsmath}\usepackage[most]{tcolorbox}\usepackage{tikz}
\usetikzlibrary{positioning,arrows.meta,fit,shapes.geometric,calc,
                backgrounds,shapes.multipart,decorations.pathreplacing}
\usepackage{adjustbox}
\usepackage{hyperref}
\geometry{top=2.4cm,bottom=2.4cm,left=2.5cm,right=2.3cm}
\setstretch{1.07}
% Laisse TeX relâcher l'espacement plutôt que de déborder dans la marge
\setlength{\emergencystretch}{1.2em}

\definecolor{mainblue}{RGB}{28,74,102}
\definecolor{darkgray}{RGB}{62,62,66}
\definecolor{lightblue}{RGB}{238,245,250}
\definecolor{ent}{RGB}{221,236,247}    % étage de la cascade
\definecolor{tab}{RGB}{215,237,221}    % donnée
\definecolor{mod}{RGB}{233,233,238}    % module logiciel
\definecolor{ext}{RGB}{250,233,207}    % acteur ou système externe
\definecolor{alerte}{RGB}{250,226,226} % risque, interdit
\definecolor{codebg}{RGB}{246,246,248}
\definecolor{okgreen}{RGB}{46,125,50}
\definecolor{kored}{RGB}{176,42,42}

\hypersetup{colorlinks=true,linkcolor=mainblue,citecolor=mainblue,urlcolor=mainblue,
  pdftitle={Assistant bancaire frugal - Proposition de solution, sujet 1},
  pdfauthor={Equipe Yahya, Sohayb, Younes}}

\titleformat{\chapter}[display]
  {\normalfont\huge\bfseries\color{mainblue}}
  {\normalsize\mdseries\color{darkgray}\MakeUppercase{\chaptertitlename}\ \thechapter}
  {6pt}{\Huge}
\titlespacing*{\chapter}{0pt}{-24pt}{22pt}
\titleformat{\section}{\Large\bfseries\color{mainblue}}{\thesection}{0.7em}{}
\titleformat{\subsection}{\large\bfseries\color{darkgray}}{\thesubsection}{0.7em}{}
\titleformat{\subsubsection}{\normalsize\bfseries\color{darkgray}}{\thesubsubsection}{0.6em}{}
\titlespacing*{\section}{0pt}{1.4em}{0.5em}
\titlespacing*{\subsection}{0pt}{1.05em}{0.35em}

\captionsetup{font=small,labelfont=bf}
\setlist[itemize]{leftmargin=1.4em,itemsep=0.12em,topsep=0.25em}
\setlist[enumerate]{leftmargin=1.7em,itemsep=0.12em,topsep=0.25em}
\newcolumntype{Y}{>{\raggedright\arraybackslash}X}
\newcolumntype{Z}{>{\raggedright\arraybackslash\small}X}
\newcolumntype{C}[1]{>{\centering\arraybackslash}p{#1}}
\newcolumntype{L}[1]{>{\raggedright\arraybackslash}p{#1}}

\pagestyle{fancy}\fancyhf{}
\lhead{\small\color{darkgray}\leftmark}\cfoot{\small\thepage}
\renewcommand{\headrulewidth}{0.3pt}
\renewcommand{\chaptermark}[1]{\markboth{\chaptername\ \thechapter\ --- #1}{}}

\newtcolorbox{choix}{breakable,colback=ext!45,colframe=darkgray,boxrule=0.5pt,arc=1.4mm,
  left=3mm,right=3mm,top=2mm,bottom=2mm}
\newtcolorbox{cle}{breakable,colback=lightblue,colframe=mainblue,boxrule=0.7pt,arc=1.4mm,
  left=3mm,right=3mm,top=2mm,bottom=2mm}
\newtcolorbox{attention}{breakable,colback=alerte!60,colframe=kored!80,boxrule=0.6pt,arc=1.4mm,
  left=3mm,right=3mm,top=2mm,bottom=2mm}
\newtcolorbox{code}{breakable,colback=codebg,colframe=darkgray!35,boxrule=0.4pt,arc=1mm,
  left=2mm,right=2mm,top=1mm,bottom=1mm}

% Identifiants (intentions, fichiers, champs) : coupure autorisée après chaque souligné.
% Le souligné est pris dans la fonte : passer par \_ redéfini bouclerait.
% Robuste, car \f apparaît dans des titres et des légendes.
\DeclareRobustCommand{\f}[1]{{\texttt{\small
  \renewcommand{\_}{\char`\_\allowbreak}#1}}}
\newcommand{\eur}{\,\texteuro}
\newcommand{\oui}{\textcolor{okgreen}{\textbf{oui}}}
\newcommand{\non}{\textcolor{kored}{\textbf{non}}}
\newcommand{\partiel}{\textcolor{darkgray}{\textbf{partiel}}}
\newcommand{\amesurer}{\textit{\textcolor{darkgray}{à mesurer}}}

% ── Styles de diagrammes, constants dans tout le rapport ──
\tikzset{
  etage/.style={rectangle,draw=mainblue,fill=ent,rounded corners=2pt,
                align=center,font=\scriptsize,inner sep=4pt,minimum height=0.8cm},
  donnee/.style={rectangle,draw=mainblue!70,fill=tab,rounded corners=2pt,
                 align=center,font=\scriptsize,inner sep=4pt,minimum height=0.7cm},
  module/.style={rectangle,draw=darkgray!70,fill=mod,rounded corners=2pt,
                 align=center,font=\scriptsize,inner sep=4pt,minimum height=0.7cm},
  acteur/.style={rectangle,draw=darkgray,fill=ext,rounded corners=2pt,
                 align=center,font=\scriptsize,inner sep=4pt,minimum height=0.7cm},
  risque/.style={rectangle,draw=kored!80,fill=alerte,rounded corners=2pt,
                 align=center,font=\scriptsize,inner sep=4pt,minimum height=0.7cm},
  test/.style={diamond,aspect=2.2,draw=mainblue,fill=white,align=center,
               font=\scriptsize,inner sep=1pt},
  fl/.style={-{Latex[length=1.8mm]},semithick,draw=mainblue!80},
  flp/.style={-{Latex[length=1.8mm]},semithick,draw=darkgray!75,dashed},
  lien/.style={semithick,draw=darkgray!70},
  et/.style={font=\scriptsize,text=darkgray,inner sep=1.5pt},
  gantt/.style={rectangle,draw=mainblue!80,rounded corners=1pt,font=\tiny,
                align=center,inner sep=1.5pt,minimum height=0.46cm},
}

\begin{document}

\begin{titlepage}
\centering
\vspace*{0.6cm}
{\large\scshape\color{darkgray} Hackathon IA\\[0.15cm]
La Banque Postale $\times$ CGI $\times$ ENSEIRB-MATMECA\par}
\vspace{0.3cm}
{\normalsize\color{darkgray} Sujet 1 --- Assistant bancaire personnalisé\par}

\vspace{2.6cm}
{\large\color{darkgray} Proposition de solution et feuille de route\par}
\vspace{0.5cm}
{\Huge\bfseries\color{mainblue} Un assistant bancaire\\[0.25cm] frugal, explicable, et qui agit\par}
\vspace{0.6cm}
{\Large Comprendre, personnaliser, conseiller\\[0.1cm]
avec le plus petit calcul qui suffit\par}

\vspace{1.8cm}
\rule{0.78\textwidth}{0.7pt}
\vspace{1.1cm}

\begin{tabular}{rl}
\textbf{Yahya Bel Hajjam} & Intelligence artificielle \\[0.15cm]
\textbf{Sohayb} & Génie logiciel \\[0.15cm]
\textbf{Younes} & Calcul intensif \\
\end{tabular}

\vfill
{\small\color{darkgray} Document de travail interne à l'équipe --- lundi 26 octobre 2026\\
Soutenance : mercredi 28 octobre, 13h30\par}
\end{titlepage}

\tableofcontents
\cleardoublepage

\chapter*{La solution en une page}
\addcontentsline{toc}{chapter}{La solution en une page}
\markboth{La solution en une page}{}

\paragraph{Le problème.} Un client écrit à sa banque en langage libre : \og Je veux acheter un
appartement \fg, \og J'ai perdu ma carte hier soir \fg, \og Pourquoi j'ai payé 8\eur{} de frais ? \fg.
L'assistant doit comprendre, poser les bonnes questions, tenir compte de la situation du client, et
proposer une piste --- avec une IA open source, locale, la plus sobre possible, et dont chaque
conseil se justifie.

\paragraph{Notre réponse : une cascade.} Chaque étage ne passe la main au suivant que s'il ne sait
pas faire lui-même. Le calcul le plus cher n'est appelé qu'en dernier recours, et c'est mesuré.

\begin{center}\small
\begin{tabularx}{\textwidth}{@{}l L{3.6cm} Y@{}}
\toprule
\textbf{Étage} & \textbf{Qui fait le travail} & \textbf{Pourquoi pas un LLM} \\
\midrule
1. Comprendre & Mots-clés, sinon petit encodeur (118 M de paramètres) et régression logistique &
  Les cas évidents coûtent moins d'une milliseconde, les autres quelques millisecondes ; tout se
  mesure sur le jeu de test officiel de Banking77. \\
2. Converser & Parcours décrits en YAML : quoi demander, dans quel ordre & Les questions à poser se
  déduisent du profil : on ne demande jamais ce qu'on sait déjà. \\
3. Conseiller & Règles publiques (plafonds, 35\,\% d'endettement) et modèle d'appétence &
  Une décision financière doit être reproductible et justifiable, pas générée. \\
4. Agir & Outils : opposition, virement guidé, transmission au conseiller & Une action ne part
  qu'après confirmation explicite du client. \\
5. Formuler & Le seul appel LLM du tour, à la fin : \f{ministral-3:3b} (3,0 Go) sous garde-fou ;
  gabarits pour les réponses simples & Le LLM met en forme des faits calculés ; il n'en invente
  aucun. \\
Escalade & \f{ministral-3:14b} (9,1 Go), seulement si l'étage 1 doute encore après une question &
  Le \og gros modèle à un instant précis \fg{} que CGI a explicitement encouragé. \\
\bottomrule
\end{tabularx}
\end{center}

\paragraph{Les quatre décisions qui font la différence.}
\begin{enumerate}
  \item \textbf{Comprendre sans LLM.} Des mots-clés pour les cas évidents, un petit classifieur
        pour le reste. Il est entraîné sur Banking77 traduit en français, plus 8 intentions que
        Banking77 ne couvre pas : épargne, crédit immobilier, crédit à la consommation, prêt
        étudiant, virement à un proche, jargon, insatisfaction, hors sujet.
  \item \textbf{Décider sans LLM.} Chaque proposition vient avec ses raisons \emph{et} avec la
        liste des produits écartés et pourquoi (\og pourquoi pas le LEP ? \fg).
  \item \textbf{Formuler avec le plus petit modèle qui suffit, une seule fois par tour.} On part
        d'un modèle qui parle bien (14B), on descend en taille, et on garde le plus petit qui tient
        encore le seuil de qualité.
  \item \textbf{Tout mesurer contre une version \og tout LLM \fg.} Même jeu de scénarios, trois
        configurations, et un tableau qualité / énergie / latence / taille au cœur du pitch.
\end{enumerate}

\paragraph{Ce qu'un organisateur nous a dit.} \og Pour gagner, c'est moins de LLM, plus d'algo, ou
des modèles d'IA plus spécialisés \fg, y compris pour le sujet 1 : des mots-clés et un classifieur
en entrée, un LLM appelé une seule fois, à la fin, et dont on choisit la taille en descendant.
C'est l'architecture ci-dessus.

\needspace{10\baselineskip}
\paragraph{Ce que la grille du jury y gagne.}
\begin{center}\small
\begin{tabularx}{\textwidth}{@{}l c Y@{}}
\toprule
\textbf{Critère} & \textbf{Poids} & \textbf{Notre preuve} \\
\midrule
Pertinence métier & 25\,\% & Les 6 situations des slides, plus l'appartement, le prêt étudiant et
  le jargon ; priorité n\textsuperscript{o}\,1 de La Banque Postale ; inclusion et accessibilité. \\
Qualité technique & 20\,\% & Précision sur le test officiel Banking77 (anglais et français) ; taux
  de réussite d'un jeu de 30 scénarios automatisés. \\
Frugalité & 20\,\% & Énergie par requête (CodeCarbon), latence, Go chargés, part des requêtes
  traitées sans LLM ; comparaison à la version \og tout LLM \fg. \\
Explicabilité & 15\,\% & Panneau \og pourquoi \fg{} : intention et confiance, règles appliquées,
  produits proposés et écartés, sources publiques. \\
Démonstration & 10\,\% & Six minutes en direct, avec les profils fournis par les slides ; vidéo de
  secours. \\
Présentation & 10\,\% & Trois voix, une par module ; un tableau de résultats, une démo, un
  message. \\
\bottomrule
\end{tabularx}
\end{center}

\paragraph{L'organisation.} Trois modules, un responsable et un suppléant chacun : \emph{Comprendre
et mesurer} (Younes), \emph{Converser et formuler} (Yahya), \emph{Conseiller et agir} (Sohayb).
Les contrats entre modules sont figés lundi soir ; deux intégrations mardi (12h30 et 18h) ; gel
fonctionnel mardi 18h ; mercredi matin consacré aux mesures finales, à la vidéo et aux répétitions.

\begin{cle}
\textbf{La phrase à retenir.} \emph{Notre assistant comprend souvent avec de simples mots-clés,
sinon avec un modèle de 118 millions de paramètres ; il décide avec des règles que le client peut
lire ; et il n'appelle un LLM qu'une fois, à la fin, pour mettre en mots ce qui a déjà été
calculé.}
\end{cle}

\chapter{Ce que le hackathon attend}
\label{ch:attentes}

Ce chapitre rassemble tout ce qui a été demandé, écrit ou dit, sans interprétation. Il se termine
par une liste d'exigences numérotées : le chapitre~\ref{ch:couverture} vérifie que la solution
répond à chacune.

\section{La règle commune aux deux sujets}

Les slides de CGI posent quatre règles, \og deux sujets, une même règle : une IA frugale et open
source pour le domaine bancaire \fg.

\begin{center}\small
\begin{tabularx}{\textwidth}{@{}L{4cm} Y@{}}
\toprule
\textbf{Règle} & \textbf{Ce qu'elle impose} \\
\midrule
Open source uniquement & Licences ouvertes, aucune API de LLM externe, données en local. \\
Le plus petit modèle qui suffit & Le GPU est partagé : une base simple d'abord, un modèle plus gros
  seulement si le gain est \emph{mesuré}. \\
Serveur GPU sur site & Données et modèles restent en interne. \\
Mesurer sa frugalité & Taille du modèle, énergie (CodeCarbon) et latence, présentées dans le
  pitch. \\
\bottomrule
\end{tabularx}
\end{center}

Denis Barthou a précisé : la solution tourne uniquement sur les machines locales ; utiliser d'autres
IA pour \emph{développer} est autorisé ; dans la solution, on privilégie de petits LLM, \og des
algorithmes au lieu de LLM quand c'est possible \fg, des bases de données et d'autres types d'IA
(classifieurs). Le critère de frugalité retient le plus petit LLM \emph{en Go, quantisation
comprise}.

À l'oral, CGI a ajouté trois nuances qui orientent directement notre architecture :
\begin{itemize}
  \item la frugalité, c'est peu de tokens, \og mais surtout \fg{} peu d'énergie ;
  \item \og le plus petit modèle qui suffit \fg{} est une phrase \og à double tranchant \fg : une
        architecture peut appeler \textbf{un gros modèle à un instant précis, pour un besoin
        précis}, et de petits modèles le reste du temps --- \og c'est là que vous allez innover \fg ;
  \item le sujet 1 est \og de l'IA générative qu'on orchestre \fg, \og du chatbot, mais pas du
        chatbot traditionnel \fg.
\end{itemize}

\section{Le sujet 1}

\begin{cle}
\og Un client écrit : \emph{Je veux acheter un appartement.} Votre assistant doit comprendre, poser
les bonnes questions et lui proposer une piste adaptée à sa situation. \fg
\end{cle}

\paragraph{Objectifs.} Comprendre les questions des clients en langage naturel ; personnaliser
selon l'âge, les revenus et l'historique ; conseiller, c'est-à-dire recommander un produit bancaire
simple.

\paragraph{Démarche attendue, en quatre temps.} \emph{Cadrer} (quels besoins et quelles intentions
prioriser ?), \emph{Préparer} (rendre les données exploitables en français), \emph{Concevoir}
(quelle architecture, et pourquoi la plus sobre ?), \emph{Démontrer} (prouver la valeur et la
frugalité).

\paragraph{Livrables attendus.} Données préparées en français ; modèle de compréhension des
intentions ; chatbot relié à une base clients simulée ; moteur de conseils personnalisés ; mesures
de frugalité.

\paragraph{Données fournies.} Banking77 (intentions) ; Bank Marketing, \og enrichi de revenus et
d'historique simulés \fg{} ; le catalogue de produits est à construire à partir d'informations
publiques.

\paragraph{Les situations données en exemple.} Six sur les slides, trois de plus à l'oral.

\begin{center}\small
\begin{tabularx}{\textwidth}{@{}l L{4.3cm} L{3.2cm} Y@{}}
\toprule
\textbf{Thème} & \textbf{Message} & \textbf{Profil} & \textbf{Comportement attendu} \\
\midrule
Carte & \og J'ai perdu ma carte hier soir, je fais quoi ? \fg & Inès, 34 ans, cadre & Réagir vite,
  guider l'opposition, rassurer \\
Réclamation & \og Pourquoi j'ai payé 8\eur{} de frais ce mois-ci ? \fg & Marc, 52 ans, découverts
  fréquents & Expliquer les frais, orienter si besoin \\
Épargne & \og Je voudrais mettre de côté pour mes enfants. \fg & Sophie, 41 ans, deux enfants &
  Demander l'horizon, proposer une piste \\
Virement & \og Je dois envoyer 500\eur{} à mon fils. \fg & Paul, 68 ans, peu à l'aise en ligne &
  Guider pas à pas, avec des mots simples \\
Crédit & \og Je veux un crédit pour acheter une voiture. \fg & Karim, 29 ans, 1\,900\eur{} de
  revenus & Questionner, rester prudent sur l'endettement \\
Insatisfaction & \og Votre appli ne marche jamais ! \fg & Client mécontent, 3\textsuperscript{e}
  message du jour & Reconnaître la frustration, passer la main \\
\midrule
Logement & \og Je veux acheter un appartement. \fg & (énoncé principal) & Comprendre, questionner,
  proposer une piste \\
Prêt étudiant & \og Je veux faire un prêt étudiant, comment je fais ? \fg & (exemple oral) &
  Regard d'un jeune client \\
Jargon & Un mot ou un sigle bancaire non compris & (exemple oral) & Expliquer avec des mots simples \\
\bottomrule
\end{tabularx}
\end{center}

\paragraph{Précisions données à l'oral.} Les jeux de données sont \og indicatifs \fg{} : on peut en
croiser d'autres. La profondeur est à notre choix : \og un agent qui va jusqu'au bout d'un workflow
et qui fait lui-même l'action \fg{} (couper la carte selon la gravité), ou un assistant qui guide
seulement. Ils veulent \og notre regard \fg{} de jeunes clients sur le niveau de service attendu. Le
chatbot doit répondre \og immédiatement \fg. Les attentes seront précisées avec chaque équipe au
fil des jours.

\section{Ce que le jury regarde}

Jury de trois personnes : Adrien (CGI, aussi coach), Denis Barthou (ENSEIRB) et Jean-Michel (La
Banque Postale, à distance ; les équipes peuvent lui poser des questions). Il \og évalue votre
démarche autant que votre résultat \fg.

\begin{center}\small
\begin{tabularx}{\textwidth}{@{}c l Y@{}}
\toprule
\textbf{Poids} & \textbf{Critère} & \textbf{Formulation officielle} \\
\midrule
25\,\% & Pertinence métier & Utile pour un client ou un analyste \\
20\,\% & Qualité technique & Résultats mesurés et bien évalués \\
20\,\% & Frugalité & Taille, énergie, latence \\
15\,\% & Explicabilité & Chaque conseil ou alerte se justifie \\
10\,\% & Démonstration & Elle fonctionne et se comprend \\
10\,\% & Présentation & Claire et tenue en 15 minutes \\
\bottomrule
\end{tabularx}
\end{center}

\paragraph{Format.} 30 minutes par équipe (20 s'il y a plus de 7 groupes) : 15 minutes de
présentation avec \textbf{démonstration en direct sur le serveur fourni}, environ 10 minutes de
questions, 5 minutes de retour à chaud. À remettre avant : le code source déposé sur un dépôt Git
et le support de présentation.

\section{Ce que La Banque Postale valorise}

La directrice SI de La Banque Postale a présenté ses priorités IA dans cet ordre :
\begin{enumerate}
  \item \textbf{l'expérience client} : chatbot conversationnel sur ses contrats, \textbf{les
        réclamations} (\og un sujet prégnant \fg), et la personnalisation des offres à partir des
        avoirs du client --- c'est exactement le sujet 1 ;
  \item l'optimisation des processus (dossiers de crédit, anti-blanchiment, anti-fraude) ;
  \item l'innovation (outils pour les collaborateurs, incubateur).
\end{enumerate}
Elle a insisté sur trois valeurs : \textbf{l'accessibilité} (clients et collaborateurs malvoyants
ou malentendants : tout ce que construit la DSI doit être compatible), \textbf{l'inclusion
bancaire} (la banque a le devoir d'accueillir tous les clients, y compris les plus fragiles) et la
\textbf{responsabilité} (banque citoyenne, finance durable). CGI a mentionné la souveraineté et
la RSE parmi les enjeux de ses clients. Un client sur quatre en France est client de La Banque
Postale, accueilli dans 17\,000 points de contact.

\section{Les contraintes}

\begin{itemize}
  \item \textbf{Le temps.} Lundi soir, mardi, mercredi matin : un jour et demi de travail effectif.
  \item \textbf{Les machines.} Postes de la salle (RTX 3060 12 Go, 30 Go de RAM) : machine
        principale de développement. DeepEirb (2 $\times$ A6000 48 Go) : SSH et serveur
        d'inférence vLLM, à condition de \emph{demander à Denis Barthou} les modèles voulus.
        Jetson AGX Orin et OrangePi : disponibles mardi, pour le déploiement uniquement.
  \item \textbf{Les modèles déjà installés.} \f{gemma4:e2b} (7,2 Go) et \f{ministral-3:14b}
        (9,1 Go) sous Ollama, contexte réglé à 16\,384 tokens. Télécharger d'autres modèles est
        possible, \og pas en wifi \fg.
  \item \textbf{L'équipe.} Trois personnes : Yahya (IA), Sohayb (génie logiciel), Younes (calcul
        intensif).
\end{itemize}

\section{Ce qu'un organisateur nous a conseillé}

Lundi vers 17h30, un organisateur du hackathon a pris le temps d'échanger avec nous
(transcription dans \f{infos/transcription-3-echange-organisateur.txt}). Ses conseils précisent
la règle commune :

\begin{itemize}
  \item \textbf{\og Pour gagner, c'est moins de LLM, plus d'algo \fg}, ou des algorithmes différents,
        ou des modèles d'IA plus spécialisés --- \og même pour le premier \fg{} sujet.
  \item \textbf{En entrée, pas besoin de LLM} pour savoir de quoi parle le client : un algorithme de
        catégorisation est plus rapide. Une simple recherche de \textbf{mots-clés}, presque gratuite,
        trouve le sujet dans environ 30\,\% des cas (\og ça fait 30\,\% de gagné \fg), et peut
        donner \textbf{un degré de confiance} : un seul mot-clé trouvé dans une longue phrase, ce
        n'est pas sûr. Les autres cas passent à un outil plus lourd.
  \item \textbf{En sortie, un LLM}, parce qu'il faut parler français à un humain ; mais il faut
        \textbf{minimiser les appels} et ne l'appeler \textbf{qu'à la fin}.
  \item \textbf{Choisir sa taille en descendant} : commencer avec un modèle qui parle correctement
        (\og un 15B ou un 30B \fg), puis descendre vers un 4B, un 2B (\og même le Gemma 4 que vous
        avez \fg), voire un 1B.
  \item \textbf{Se demander, à chaque endroit, si un LLM est nécessaire} : jamais, ou seulement pour
        une partie des cas. \og Seulement 30\,\% des cas avec un LLM, c'est intéressant \fg{} : on
        prend le vélo pour aller chercher le pain, la voiture pour le week-end.
  \item \textbf{Un seul sujet}, bien fait. Puis, \textbf{une fois la solution sûre}, essayer des
        idées nouvelles, même exotiques ; utiliser l'IA pour coder et discuter des idées ; et
        \og se faire plaisir \fg.
\end{itemize}

\section{Les exigences, numérotées}
\label{sec:exigences}

\begin{center}\small
\begin{longtable}{@{}l L{8.6cm} L{4.4cm}@{}}
\toprule
\textbf{N\textsuperscript{o}} & \textbf{Exigence} & \textbf{Source} \\
\midrule
\endhead
E1 & Comprendre les demandes en langage naturel, en français & Objectif 01, livrable \\
E2 & Fournir des données préparées en français & Livrable \\
E3 & Poser les bonnes questions avant de proposer & Énoncé \\
E4 & Personnaliser selon l'âge, les revenus et l'historique & Objectif 02 \\
E5 & Recommander un produit bancaire simple & Objectif 03, livrable \\
E6 & Relier le chatbot à une base clients simulée & Livrable \\
E7 & Construire le catalogue à partir d'informations publiques & Données \\
E8 & Open source, aucune API externe, données locales & Règle commune \\
E9 & Le plus petit modèle qui suffit ; un plus gros seulement si le gain est mesuré & Règle commune \\
E10 & Mesurer taille, énergie (CodeCarbon) et latence & Règle commune, livrable \\
E11 & Justifier chaque conseil & Grille (15\,\%) \\
E12 & Répondre immédiatement & Oral \\
E13 & Choisir et assumer la profondeur : guider ou agir & Oral \\
E14 & Traiter les six situations avec le comportement attendu & Slide 3 \\
E15 & Expliquer le jargon ; apporter le regard d'un jeune client & Oral \\
E16 & Démo en direct sur le serveur fourni ; code sur Git ; support remis & Soutenance \\
E17 & Respecter accessibilité et inclusion (souhaitable) & Présentation de La Banque Postale \\
E18 & En entrée, des algorithmes ; un seul appel au LLM, à la fin ; mesurer la part des cas sans
  LLM & Échange avec un organisateur \\
E19 & D'abord une solution sûre, ensuite des idées nouvelles & Échange avec un organisateur \\
\bottomrule
\end{longtable}
\end{center}

\chapter{Ce que disent les données}
\label{ch:donnees}

Les deux jeux de données ont été téléchargés et lus le 26 octobre. Les chiffres de ce chapitre sont
comptés sur les fichiers, pas recopiés d'une documentation.

\section{Banking77 : la compréhension des demandes}

\paragraph{Ce que c'est.} 13\,083 questions de clients d'une banque en ligne britannique, chacune
étiquetée par l'une de 77 intentions (Casanueva et al., 2020, licence CC-BY-4.0). Le jeu
d'entraînement compte 10\,003 questions, réparties de façon inégale (de 35 à 187 par intention) ;
le jeu de test en compte 3\,080, exactement 40 par intention. \textbf{Tout est en anglais.}

\paragraph{Les scores de référence.} L'article d'origine donne, en précision sur le test complet :
87,19\,\% pour BERT utilisé comme encodeur figé, 92,81\,\% pour USE, 93,01\,\% pour ConveRT,
93,36\,\% pour USE et ConveRT combinés, 93,66\,\% pour BERT affiné. Le message utile pour nous :
\textbf{un encodeur figé et un classifieur léger se placent à un ou deux points du meilleur modèle
affiné}, pour un coût d'entraînement de quelques minutes sur CPU.

\paragraph{Ce qu'il couvre, et surtout ce qu'il ne couvre pas.} Le vocabulaire est celui d'une
néobanque : recharges (\emph{top-up}), cartes virtuelles, Apple Pay, taux de change. Confronté aux
situations du sujet, le constat est net :

\begin{center}\small
\begin{tabularx}{\textwidth}{@{}l c Y@{}}
\toprule
\textbf{Situation du sujet} & \textbf{Couverte ?} & \textbf{Intentions Banking77 concernées} \\
\midrule
Carte perdue (Inès) & \oui & \f{lost\_or\_stolen\_card}, \f{compromised\_card},
  \f{card\_swallowed} \\
Frais (Marc) & \oui & \f{extra\_charge\_on\_statement}, \f{card\_payment\_fee\_charged},
  \f{transfer\_fee\_charged}, \f{cash\_withdrawal\_charge}\dots \\
Virement à un proche (Paul) & \partiel & Les intentions \f{transfer\_*} portent sur l'\emph{état}
  d'un virement (en attente, refusé, annulé), pas sur \og comment en faire un \fg \\
Épargne (Sophie) & \non & --- \\
Crédit auto (Karim) & \non & --- \\
Appartement (énoncé) & \non & --- \\
Prêt étudiant (oral) & \non & --- \\
Insatisfaction & \non & --- \\
Jargon (oral) & \non & --- \\
\bottomrule
\end{tabularx}
\end{center}

\begin{cle}
\textbf{Banking77 couvre le \emph{service}, pas le \emph{conseil}.} Il sait reconnaître une carte
perdue ou un frais contesté ; il ne contient pas une seule demande d'épargne, de crédit ou de projet
immobilier --- c'est-à-dire l'exemple même de l'énoncé. Les intentions de conseil sont à créer, en
français, par nous. C'est une partie du livrable \og données préparées en français \fg.
\end{cle}

\paragraph{Un bruit d'étiquetage à connaître.} Certaines étiquettes sont discutables. Dans le jeu
d'entraînement, \og \emph{What is the word?} \fg{} est rangé sous
\f{direct\_debit\_payment\_not\_recognised}, et \og \emph{Has there been any activity on my card
today?} \fg{} sous \f{lost\_or\_stolen\_card}. Cela plafonne la précision atteignable et justifie de
raisonner aussi au niveau des \emph{familles} d'intentions, plus robustes (section~\ref{sec:nlu}).

\paragraph{Neuf familles.} Nous regroupons les 77 intentions en neuf familles, vérifié par script :
chaque intention est rangée une fois et une seule. Le détail est en annexe~\ref{an:taxonomie}.

\begin{center}\small
\begin{tabular}{@{}l r r r@{}}
\toprule
\textbf{Famille} & \textbf{Intentions} & \textbf{Entraînement} & \textbf{Test} \\
\midrule
\f{carte\_perdue} & 4 & 350 & 160 \\
\f{carte\_usage} & 10 & 1\,084 & 400 \\
\f{carte\_commande} & 12 & 1\,465 & 480 \\
\f{frais\_contestation} & 14 & 2\,225 & 560 \\
\f{virement} & 10 & 1\,409 & 400 \\
\f{compte\_identite} & 8 & 884 & 320 \\
\f{especes} & 5 & 764 & 200 \\
\f{change\_etranger} & 6 & 815 & 240 \\
\f{alimenter\_compte} & 8 & 1\,007 & 320 \\
\midrule
\textbf{Total} & \textbf{77} & \textbf{10\,003} & \textbf{3\,080} \\
\bottomrule
\end{tabular}
\end{center}

\section{Bank Marketing : la base clients}

\paragraph{Ce que c'est.} Les clients d'une banque portugaise contactés par téléphone entre mai 2008
et novembre 2010 pour souscrire un dépôt à terme (UCI, Moro et al.). Le fichier
\f{bank-full.csv} compte \textbf{45\,211 clients} et 17 colonnes ; la variante
\f{bank-additional-full.csv} (41\,188 lignes) remplace le solde par des indicateurs économiques.
Nous partons de \f{bank-full.csv}, qui contient le solde.

\begin{center}\small
\begin{tabularx}{\textwidth}{@{}L{3.3cm} Y L{4.5cm}@{}}
\toprule
\textbf{Colonnes} & \textbf{Contenu} & \textbf{Usage pour nous} \\
\midrule
\f{age}, \f{job}, \f{marital}, \f{education} & Âge (18 à 95 ans, médiane 39), 12 métiers dont 938
  étudiants et 2\,264 retraités, situation (\og divorced \fg{} inclut les veufs), études &
  Profil, éligibilité (Livret Jeune, LDDS), registre de langue \\
\f{balance} & Solde annuel moyen en euros, de $-$8\,019 à 102\,127 ; médiane 448 ;
  \textbf{3\,766 soldes négatifs (8,3\,\%)} & Découverts fréquents (Marc), capacité d'épargne \\
\f{housing}, \f{loan}, \f{default} & Crédit immobilier en cours (55,6\,\%), crédit personnel
  (16,0\,\%), défaut de paiement (1,8\,\%) & Charges existantes, prudence crédit \\
\f{contact}, \f{day}, \f{month}, \f{duration}, \f{campaign}, \f{pdays}, \f{previous},
  \f{poutcome} & Historique des campagnes de démarchage & Seulement \f{previous} et
  \f{poutcome}, pour l'appétence \\
\f{y} & A souscrit un dépôt à terme : 5\,289 oui (11,7\,\%) & Cible du modèle d'appétence épargne \\
\bottomrule
\end{tabularx}
\end{center}

\begin{attention}
\textbf{La colonne \f{duration} est une fuite.} La documentation de l'UCI le dit explicitement : la
durée de l'appel n'est connue qu'\emph{après} l'appel, et elle détermine presque la réponse (un
appel de 0 seconde donne toujours \og non \fg). Un modèle qui l'utilise affiche un excellent score
et ne sert à rien. Nous l'excluons, ainsi que \f{contact}, \f{day} et \f{month}, propres à la
campagne. Le dire au jury est un point de qualité technique gratuit.
\end{attention}

\section{Ce qui manque, et comment on le comble}

Bank Marketing ne contient ni revenus, ni enfants, ni charges, ni historique d'opérations. Or
l'énoncé demande de personnaliser \og selon l'âge, les revenus et l'historique \fg. Les slides
annoncent une version \og enrichie de revenus et d'historique simulés \fg{} fournie par les
organisateurs. Deux cas :

\begin{itemize}
  \item \textbf{Elle arrive à temps.} On la branche derrière un \emph{adaptateur} unique,
        \f{src/conseiller/clients.py}, qui expose toujours le même profil au reste du code.
  \item \textbf{Elle tarde.} On simule nous-mêmes, avec une graine fixe et une méthode écrite :
        revenu mensuel net tiré par métier et âge ; nombre d'enfants à partir de l'âge et de la
        situation ; charges de crédit si \f{housing} ou \f{loan} ; douze mois d'événements (frais,
        découverts, contacts avec le service client), plus nombreux si le solde est négatif.
\end{itemize}

\begin{choix}
\textbf{Règle de transparence.} Toute donnée simulée porte le drapeau \f{simule=True} et sa méthode
est décrite dans la fiche de données. Dans l'interface, un champ simulé s'affiche en italique. Le
jury ne doit jamais pouvoir croire qu'un revenu inventé est une donnée réelle.
\end{choix}

\paragraph{Les clients de démonstration.} Les profils des slides deviennent huit clients fixes,
chacun ancré sur une ligne réelle de Bank Marketing (âge, métier, situation, solde, crédits) et
complété par simulation. Ils servent à la démo et aux tests.

\begin{center}\small
\begin{tabularx}{\textwidth}{@{}l l L{3.5cm} Y@{}}
\toprule
\textbf{Id} & \textbf{Client} & \textbf{Ancrage Bank Marketing} & \textbf{Enrichissement simulé
  (exemple)} \\
\midrule
C01 & Inès, 34 ans & \f{management}, \f{single} & 3\,200\eur{} nets, carte active \\
C02 & Marc, 52 ans & \f{blue-collar}, \f{married}, solde négatif & 1\,850\eur{} nets, 5 découverts
  en 12 mois, commissions facturées \\
C03 & Sophie, 41 ans & \f{admin.}, \f{married} & 2\,400\eur{} nets, enfants de 8 et 12 ans,
  Livret A \\
C04 & Paul, 68 ans & \f{retired}, \f{divorced} & 1\,600\eur{}, \og peu à l'aise en ligne \fg, fils
  enregistré comme bénéficiaire \\
C05 & Karim, 29 ans & \f{technician}, \f{loan=yes} & 1\,900\eur{} nets (donné), crédit en cours
  de 150\eur{}/mois \\
C06 & Julie, 37 ans & \f{services} & 3 contacts au service client aujourd'hui \\
C07 & Léa, 31 ans & \f{management}, \f{housing=no} & 2\,900\eur{} nets, 25\,000\eur{} d'épargne \\
C08 & Lucas, 20 ans & \f{student} & Job étudiant à 450\eur{}, sans épargne \\
\bottomrule
\end{tabularx}
\end{center}

\section{Le catalogue de produits}

L'énoncé le laisse à construire \og à partir d'informations publiques \fg. Nous le limitons à ce
qu'un assistant peut honnêtement proposer --- des produits \emph{simples} --- et nous l'écrivons en
JSON, une fiche par produit : identifiant, nom, catégorie, conditions d'éligibilité, plafond,
disponibilité de l'argent, à qui il convient, limites, \textbf{source publique et date de
relevé}. La liste complète et les paramètres sont en annexe~\ref{an:catalogue}.

\begin{attention}
\textbf{Les taux réglementés changent} (le taux du Livret A est révisé deux fois par an). Aucun taux
n'est écrit en dur dans le code : c'est un champ du catalogue, avec sa date et sa source, relevé le
jour même sur le site public. Les plafonds (22\,950\eur{} pour le Livret A, 12\,000\eur{} pour le
LDDS, 10\,000\eur{} pour le LEP, 1\,600\eur{} pour le Livret Jeune, 61\,200\eur{} pour le PEL) sont
stables mais sont aussi vérifiés à la source.
\end{attention}

\chapter{Ce que Timso AI nous apprend}
\label{ch:timso}

Timso AI est l'assistant de rédaction construit pendant le stage de Yahya pour l'agence de voyage
French Touch Travel, et mis en production. Il lit le patrimoine de l'agence --- 68\,648 courriels,
15\,766 journées de voyage vendues, 4\,282 fiches de catalogue --- et s'en sert pour rédiger des
brouillons qu'un conseiller relit avant envoi.

Ce n'est pas le même problème, mais c'est la même famille : \textbf{comprendre une demande libre,
retrouver le bon produit dans un catalogue, et rédiger sans rien inventer}. Plusieurs décisions y
ont été \emph{mesurées} en conditions réelles. Les reprendre nous évite de les redécouvrir en deux
jours.

\section{Les leçons transposables}

\begin{center}\small
\begin{longtable}{@{}L{3.4cm} L{5.4cm} L{5.6cm}@{}}
\toprule
\textbf{Leçon de Timso AI} & \textbf{Ce qui a été décidé ou mesuré} & \textbf{Ce qu'on en fait ici} \\
\midrule
\endhead
Recherche plutôt qu'affinage & L'affinage d'un LLM a été écarté avant tout travail : coûteux,
  figé sur une version, il demande des GPU. & \textbf{Aucun affinage de LLM.} On n'entraîne qu'une
  tête de classification, en quelques secondes sur CPU. \\
\addlinespace
Encodeurs E5 en local & \f{multilingual-e5-large} calculé sur place : aucune donnée client ne
  sort. Oublier le préfixe \f{query:}/\f{passage:} dégrade la pertinence \emph{sans aucune
  erreur}. & Même famille en version \f{small}. Les préfixes sont posés par une seule fonction,
  reprise de \f{src/search/embed/model.py}. \\
\addlinespace
Pas de base vectorielle & 119\,107 vecteurs tenaient en 465 Mo ; un produit matrice-vecteur numpy
  répondait en une centaine de millisecondes, sans index à maintenir. & Le glossaire et le
  catalogue font quelques centaines d'entrées : \textbf{numpy en mémoire}, pas de base
  vectorielle. \\
\addlinespace
Seuils obtenus par balayage & Seuils de 0,90 et 0,88 choisis en mesurant précision et couverture.
  Une variante gagnait 2,8 points de couverture pour 0,5 de précision : \textbf{rejetée}, car une
  erreur coûte plus qu'un manque. & Le seuil de confiance du classifieur est fixé par balayage.
  En dessous, l'assistant \textbf{demande} plutôt que de deviner. \\
\addlinespace
Cascade de résolution & Chaque ligne est d'abord rapprochée du dossier en cours, puis du catalogue
  complet, et n'est liée à un produit que si elle passe le seuil de l'étage. & Même principe pour
  les modèles : petit d'abord, gros seulement si le petit doute. \\
\addlinespace
Le brouillon, jamais l'envoi & Rien ne part chez un client sans qu'un humain l'ait relu. &
  L'assistant \textbf{oriente}, le conseiller \textbf{décide} (crédit). Le passage de main
  transmet un résumé. \\
\addlinespace
Déclarer les outils, puis les exécuter & Les outils MCP séparent la déclaration (un contrat
  lisible) de l'exécution ; la lecture est libre, les effets de bord passent par un point
  d'entrée suivi. & Lecture du profil libre ; opposition et virement \textbf{seulement après
  confirmation}, avec une trace dans un journal. \\
\addlinespace
Désactiver la réflexion & Mesuré : 8\,000 tokens consommés en réflexion sans une ligne de réponse ;
  1\,571 tokens pour la réponse complète sans réflexion. & Sur les petits modèles qui raisonnent
  (Qwen3), \f{think: false}. C'est de la frugalité directe. \\
\addlinespace
Vérifier le service réel & Un indicateur de santé répondait \og vrai \fg{} parce qu'une clé était
  configurée, alors que le service était arrêté depuis quatre jours. & Avant la démo, on
  interroge vraiment Ollama. S'il ne répond pas, on bascule sur les gabarits. \\
\addlinespace
Ne pas rendre une note obligatoire & 20 notes exigées par le formulaire : seules les 6 premières
  étaient réfléchies, les 14 autres valaient 5/10 par défaut. & Grille d'évaluation humaine courte,
  champs facultatifs. \\
\bottomrule
\end{longtable}
\end{center}

\section{Ce qu'on ne reprend pas}

\begin{itemize}
  \item \textbf{L'API Claude.} Interdite ici : la règle impose des modèles open source, en local.
        C'est précisément ce qui rend l'architecture en cascade nécessaire : on ne peut plus
        compter sur un très gros modèle pour tout faire.
  \item \textbf{PostgreSQL, Docker, Keycloak.} Justifiés pour un service en production, inutiles
        pour un démonstrateur de deux jours. Des fichiers CSV et JSON suffisent.
  \item \textbf{\f{multilingual-e5-large}.} Avec 560 millions de paramètres, il est près de cinq fois
        plus gros que la version \f{small} (118 millions). On ne monte en taille que si le gain est mesuré.
\end{itemize}

\begin{cle}
\textbf{Ce que Timso AI donne au pitch.} Au moment des questions, on peut répondre par
l'expérience : \og Dans un système en production, on a mesuré qu'une erreur de rattachement coûtait
plus cher qu'un manque ; c'est pour ça que notre assistant pose une question plutôt que de
deviner. \fg{} Peu d'équipes pourront s'appuyer sur un système qui a déjà servi de vrais
utilisateurs.
\end{cle}

\chapter{La solution}
\label{ch:solution}

\section{Le principe : une cascade}

Un seul principe gouverne l'architecture : \textbf{chaque étage ne passe la main au suivant que
s'il ne sait pas faire lui-même}. Comprendre une demande ne demande pas un LLM ; décider d'un
produit ne doit pas en demander un ; mettre en mots un résultat calculé peut en demander un petit ;
lever un vrai doute peut justifier un gros, une fois. C'est la lecture littérale de \og le plus
petit modèle qui suffit \fg, et c'est aussi la réponse à la nuance de CGI sur le \og gros modèle à
un instant précis \fg. Un organisateur l'a résumé ainsi : \og pour gagner, c'est moins de LLM,
plus d'algo \fg. Dans le cas normal, un tour de conversation fait \textbf{au plus un appel à un
LLM, à la fin}.

\begin{figure}[H]
\centering
\begin{adjustbox}{max width=\textwidth}
\begin{tikzpicture}[x=1cm,y=1cm]
  % Acteurs
  \node[acteur,minimum width=3.6cm] (client) at (0,0) {Client, dans l'interface};
  \node[acteur,minimum width=5.4cm,text width=5.2cm] (out) at (0,-9.6)
    {Réponse\\+ panneau \og pourquoi \fg{} + panneau \og frugalité \fg};
  % Étages
  \node[etage,minimum width=4.4cm] (e1) at (0,-1.5)
    {\textbf{1. Comprendre}\\ mots-clés, sinon \f{multilingual-e5-small}\\ + régression logistique ; entités par règles};
  \node[test] (d) at (0,-3.05) {confiance $\geq \tau$ ?};
  \node[etage,minimum width=4.4cm] (e2) at (0,-4.65)
    {\textbf{2. Converser}\\ parcours YAML, profil du client};
  \node[etage,minimum width=3.3cm] (e3) at (-2.55,-6.35)
    {\textbf{3. Conseiller}\\ règles publiques\\ + appétence};
  \node[etage,minimum width=3.3cm] (e4) at (2.55,-6.35)
    {\textbf{4. Agir}\\ outils\\ + confirmation};
  \node[etage,minimum width=4.4cm] (e5) at (0,-8.05)
    {\textbf{5. Formuler}\\ gabarit, sinon \f{ministral-3:3b}\\ + garde-fou};
  % Doute
  \node[module,minimum width=2.9cm] (clar) at (-5.4,-3.05) {Question de\\ clarification};
  \node[risque,minimum width=2.9cm] (esc) at (-5.4,-4.65) {\textbf{Escalade}\\ \f{ministral-3:14b}};
  % Données
  \node[donnee,minimum width=2.6cm] (prof) at (-5.6,-6.35) {Profils clients\\ (simulés)};
  \node[donnee,minimum width=2.6cm] (cata) at (-5.6,-7.6) {Catalogue\\ + sources};
  \node[donnee,minimum width=2.6cm] (jour) at (5.6,-6.35) {Journal\\ des actions};
  \node[donnee,minimum width=2.6cm] (glos) at (5.6,-8.05) {Glossaire\\ (fiches vérifiées)};
  % Hors ligne
  \node[module,dashed,minimum width=3.2cm,text width=3.1cm] (off) at (5.95,-1.5)
    {\textit{Hors ligne, une fois}\\ traduction et exemples\\ en français, qui\\ entraînent l'étage 1\\ (gros modèle, DeepEirb)};
  % Mesure
  \node[module,minimum height=9.6cm,minimum width=0.75cm] (mes) at (8.15,-4.8) {};
  \node[rotate=90,font=\scriptsize,align=center] at (mes.center)
    {\textbf{Mesure continue} : latence, tokens, énergie, Go chargés};
  % Flux principal
  \draw[fl] (client) -- (e1);
  \draw[fl] (e1) -- (d);
  \draw[fl] (d) -- node[et,right]{oui} (e2);
  \draw[fl] (d) -- node[et,above]{non} (clar);
  \draw[flp] (clar.north) |- node[et,pos=0.25,left]{question} (client.west);
  \draw[flp] (clar) -- node[et,right]{si toujours ambigu} (esc);
  \draw[fl] (esc) -- (e2);
  \draw[fl] (e2) -- (e3);
  \draw[fl] (e2) -- (e4);
  \draw[fl] (e2) -- node[et,fill=white]{service simple} (e5);
  \draw[fl] (e3) -- (e5);
  \draw[fl] (e4) -- (e5);
  \draw[fl] (e5) -- (out);
  % Liens aux données
  \draw[lien] (prof) -- (e3);
  \draw[lien] (cata) -- (e3.south west);
  \draw[lien] (e4) -- (jour);
  \draw[lien] (e5) -- (glos);
  \draw[flp] (off.west) -- (e1.east);
\end{tikzpicture}
\end{adjustbox}
\caption{La cascade. Les étages en bleu fonctionnent sans gros modèle ; l'escalade (en rouge) est
l'exception, comptée et affichée.}
\label{fig:cascade}
\end{figure}

Le tableau ci-dessous fixe pour chaque étage ce qu'il reçoit, ce qu'il rend et ce qu'il coûte au
plus. Les budgets de latence sont des \emph{cibles} sur un poste de la salle (RTX 3060) ; les
valeurs réelles sont mesurées et reportées dans le pitch.

\begin{center}\small
\begin{tabularx}{\textwidth}{@{}l Y Y l@{}}
\toprule
\textbf{Étage} & \textbf{Reçoit} & \textbf{Rend} & \textbf{Budget cible} \\
\midrule
1. Comprendre & Le message, l'état du parcours en cours & Famille, intention, confiance, entités &
  $<$ 1 ms à $<$ 30 ms \\
2. Converser & La compréhension, le profil & La prochaine question, ou une demande de conseil ou
  d'action & $<$ 5 ms \\
3. Conseiller & Le besoin complété, le profil, le catalogue & Propositions, raisons, produits
  écartés, alertes & $<$ 20 ms \\
4. Agir & Une action confirmée & Un résultat (référence, ticket) & $<$ 20 ms \\
5. Formuler & Les faits structurés & Le texte affiché & 0 ms (gabarit) à $\sim$1,5 s (3B) \\
Escalade & Le message ambigu & Une intention & $\sim$2 à 4 s, rare \\
\bottomrule
\end{tabularx}
\end{center}

\section{Étage 1 --- Comprendre}
\label{sec:nlu}

\paragraph{Premier filtre : les mots-clés.} Avant tout modèle, un lexique d'expressions par famille
(\og perdu ma carte \fg, \og opposition \fg, \og frais \fg, \og virement \fg, \og mettre de côté \fg,
\og appartement \fg, \og prêt étudiant \fg, \og c'est quoi \fg\dots) est cherché dans le message
normalisé (minuscules, accents retirés, formes simples). Coût : moins d'une milliseconde, aucun
modèle chargé. Le lexique part des mots les plus discriminants de chaque famille dans les données
d'entraînement, puis il est relu à la main. Le filtre \textbf{ne conclut que s'il est sûr} : une
seule famille touchée, et soit une phrase courte, soit plusieurs expressions trouvées. Un seul
mot-clé perdu dans une longue phrase ne suffit pas : le message passe au classifieur. Une
expression n'est gardée dans le lexique que si sa précision mesurée sur la validation dépasse
97\,\%. C'est la suggestion de l'organisateur, et elle se mesure : la part des messages que les
mots-clés traitent seuls, et la précision sur cette part.

\paragraph{Classer en deux niveaux.} Le classifieur prédit d'abord une \emph{famille} (14 : les 9
familles Banking77 plus 5 nouvelles), puis l'\emph{intention} précise (85 : les 77 de Banking77
plus 8 nouvelles). La famille suffit à choisir le parcours ; l'intention sert à affiner la réponse.
Ce double niveau absorbe le bruit d'étiquetage de Banking77 : deux intentions voisines mal séparées
appartiennent presque toujours à la même famille.

\paragraph{Le modèle.} Les phrases sont transformées en vecteurs par
\f{intfloat/multilingual-e5-small} (118 millions de paramètres, vecteurs de 384 nombres, environ
0,47 Go en float32), avec le préfixe \f{query:} attendu par la famille E5. Une régression
logistique (scikit-learn) est entraînée sur ces vecteurs : quelques secondes sur CPU, un fichier de
moins d'un mégaoctet. L'encodeur étant multilingue, on l'entraîne sur trois sources à la fois :
Banking77 en anglais, Banking77 traduit en français, et nos exemples français des nouvelles
intentions. Pour expliquer ce qu'il a compris, l'assistant affiche aussi la phrase
d'entraînement la plus proche (\og votre message ressemble à : \emph{j'ai perdu ma carte} \fg),
retrouvée par un produit matrice-vecteur en mémoire : l'idée du RAG, appliquée à la compréhension.

\paragraph{Les nouvelles intentions.}

\begin{center}\small
\begin{tabularx}{\textwidth}{@{}l l Y@{}}
\toprule
\textbf{Intention} & \textbf{Famille} & \textbf{Exemple de message} \\
\midrule
\f{epargne\_projet} & \f{epargne} & \og Je voudrais mettre de côté pour mes enfants \fg \\
\f{credit\_immobilier} & \f{credit} & \og Je veux acheter un appartement \fg \\
\f{credit\_conso} & \f{credit} & \og Il me faudrait un crédit pour une voiture \fg \\
\f{pret\_etudiant} & \f{credit} & \og Comment on fait un prêt étudiant ? \fg \\
\f{faire\_virement} & \f{virement} & \og Je dois envoyer 500\eur{} à mon fils \fg \\
\f{question\_definition} & \f{aide} & \og C'est quoi un TAEG ? \fg \\
\f{insatisfaction} & \f{insatisfaction} & \og Votre appli ne marche jamais ! \fg \\
\f{hors\_sujet} & \f{conversation} & \og Bonjour \fg, \og merci \fg, \og quel temps fait-il ? \fg \\
\bottomrule
\end{tabularx}
\end{center}

Pour chacune, environ 80 exemples d'entraînement sont \emph{générés} une fois, hors ligne, par un
gros modèle, avec des consignes de variété (registre jeune ou âgé, fautes de frappe, style SMS,
phrases longues) ; un membre de l'équipe les relit tous (une vingtaine de minutes). Les 20 exemples
de \emph{test} de chaque intention sont \emph{écrits à la main} par l'équipe : tester un
classifieur sur des phrases produites par le même générateur que son entraînement surestimerait
son score.

\paragraph{Le seuil de confiance $\tau$.} Si la probabilité de la famille prédite est au moins
$\tau$, on continue. Sinon, l'assistant pose une question fermée construite à partir des deux
familles les plus probables : \og Vous parlez de votre carte, ou d'un virement ? \fg. Si la
réponse reste ambiguë, et seulement alors, le message part au modèle d'escalade. $\tau$ est choisi
par balayage sur un jeu de validation, comme dans Timso AI : on retient la plus petite valeur pour
laquelle la précision sur les messages acceptés dépasse 97\,\%. Une erreur de parcours coûte plus
qu'une question.

\paragraph{Le contexte de la conversation.} Quand un parcours attend une information (un montant,
une durée), le message est d'abord lu comme une réponse à cette question : \og 250\,000 \fg{} n'est
pas une nouvelle demande. Le classifieur n'est consulté que si le message ne répond pas à la
question en cours.

\paragraph{Les entités, par règles.} Montants (\og 500\eur{} \fg, \og 1\,900 euros \fg, \og 250k
\fg), durées (\og 10 ans \fg, \og 48 mois \fg), liens (\og mon fils \fg), objets (voiture,
appartement, études), dates relatives (\og hier soir \fg) : des expressions régulières suffisent et
coûtent zéro. Si une information manque après les règles, on ne la devine pas avec un LLM : le
parcours la demande au client. Un appel au LLM de moins, et aucun risque d'erreur d'extraction.

\paragraph{L'insatisfaction.} Deux signaux combinés : le classifieur (intention
\f{insatisfaction}) et l'historique (nombre de contacts du jour). Au troisième contact, la règle de
passage de main s'applique même si le ton est neutre.

\section{Étage 2 --- Converser}

\paragraph{Un parcours par intention de conseil ou d'action.} Chaque parcours est un fichier YAML
qui déclare les informations nécessaires, la question à poser pour chacune, l'ordre, et ce qui
se passe à la fin (un conseil, une action, un passage de main). Le gestionnaire de dialogue est un
automate qui lit ces fichiers : déterministe, testable, et modifiable sans toucher au code.

\begin{code}
\footnotesize
\begin{verbatim}
# src/converser/parcours/credit_immobilier.yaml (extrait)
famille: credit
infos:
  revenus_nets:   {source: profil, sinon: "Quels sont vos revenus nets mensuels ?"}
  credits_cours:  {source: profil, sinon: "Avez-vous des credits en cours ?"}
  prix_bien:      {question: "Quel est le prix du bien que vous visez ?"}
  apport:         {source: profil.epargne, confirmer: true,
                   question: "Quelle part de votre epargne souhaitez-vous y consacrer ?"}
fin: {conseil: credit_immobilier, mention_legale: credit, proposer: rdv_conseiller}
\end{verbatim}
\end{code}

\paragraph{Ne jamais demander ce qu'on sait.} C'est la forme la plus visible de la
personnalisation (E4) : si le profil connaît les revenus, l'assistant \emph{confirme} au lieu de
demander (\og J'ai noté vos revenus de 2\,900\eur{} nets par mois \fg). Une information simulée
est confirmée, jamais affirmée.

\paragraph{Le registre.} Un mode \og mots simples \fg{} s'active de lui-même pour un client âgé de
70 ans ou plus ou signalé \og peu à l'aise en ligne \fg{} (Paul), et à la demande pour tout le
monde : phrases courtes, une seule étape par message, aucun sigle sans explication. Ce n'est pas un
autre modèle, c'est un autre jeu de gabarits et une consigne de plus au petit LLM.

\paragraph{Le jargon.} Le glossaire (annexe~\ref{an:glossaire}) est une liste de fiches courtes,
rédigées et vérifiées par nous. Quand un terme du glossaire apparaît dans une réponse, il est
souligné et explicable d'un clic ; quand le client demande \og c'est quoi \dots{} \fg, la fiche est
retrouvée par similarité de vecteurs, en mémoire, sans LLM.

\section{Étage 3 --- Conseiller}

\paragraph{Trois passes.} Le moteur de conseil applique d'abord les conditions
\textbf{d'éligibilité} (dures : âge, revenu, plafond déjà atteint), puis des règles de
\textbf{pertinence} (horizon, objectif, épargne de précaution existante), puis un
\textbf{classement}. Chaque règle a un identifiant et une phrase d'explication paramétrée : quand
elle s'applique, la phrase remplie entre dans les raisons de la proposition.

\paragraph{Pourquoi, et pourquoi pas.} La sortie contient les produits proposés avec leurs raisons
\emph{et} les produits écartés avec la règle qui les écarte. C'est l'explication la plus
convaincante qui soit pour un client --- \og le LEP n'est pas proposé parce que votre revenu fiscal
dépasse le plafond \fg{} --- et pour le jury (15\,\% de la note).

\paragraph{L'appétence pour l'épargne.} Bank Marketing dit, pour 45\,211 clients, qui a souscrit un
dépôt à terme. Une régression logistique entraînée sur les seules caractéristiques du client (âge,
métier, situation, études, solde, crédits, campagnes précédentes ; \textbf{sans} \f{duration})
donne une probabilité d'intérêt pour l'épargne. Elle sert à \emph{ordonner} les propositions
d'épargne et à expliquer cet ordre par les trois facteurs qui pèsent le plus. On mesure l'aire sous
la courbe ROC sur un jeu de test tenu à part ; LightGBM ne remplace la régression logistique que si
le gain est mesuré. Les limites sont dites : données portugaises de 2008 à 2010, donc un signal
faible et illustratif, jamais une décision.

\paragraph{La prudence sur le crédit.} Pour un crédit immobilier, les règles sont celles du Haut
Conseil de stabilité financière : \textbf{mensualités au plus égales à 35\,\% des revenus nets,
assurance comprise, et durée au plus de 25 ans}. La capacité d'emprunt se calcule à partir de la
mensualité maximale $m$, du taux mensuel $i$ et du nombre de mois $n$ :
\[
  m = 0{,}35 \times R - C, \qquad
  \text{capacité} = m \times \frac{1 - (1+i)^{-n}}{i}
\]
où $R$ est le revenu net mensuel et $C$ les mensualités déjà en cours. Le taux est un
\emph{paramètre daté} du catalogue, affiché comme \og taux indicatif \fg. Pour le crédit à la
consommation, la même logique de taux d'effort s'applique, avec la mention légale obligatoire :
\og Un crédit vous engage et doit être remboursé. Vérifiez vos capacités de remboursement avant de
vous engager. \fg{} L'assistant ne promet jamais un accord : il donne une estimation et propose un
rendez-vous avec un conseiller, qui décide.

\section{Étage 4 --- Agir}

L'oral laissait le choix entre guider et agir (E13). Nous agissons, mais avec une règle reprise de
Timso AI : \textbf{la lecture est libre, l'effet de bord demande une confirmation explicite}, et
chaque action laisse une trace dans un journal.

\begin{center}\small
\begin{tabularx}{\textwidth}{@{}l Y c Y@{}}
\toprule
\textbf{Outil} & \textbf{Ce qu'il fait (sur le SI simulé)} & \textbf{Confirmation} &
  \textbf{Rend} \\
\midrule
\f{consulter\_profil} & Lit le profil et l'historique du client & non & Le profil \\
\f{expliquer\_frais} & Retrouve les frais du mois dans l'historique et leur cause & non & Liste
  datée des frais, avec la règle qui les explique \\
\f{bloquer\_carte} & Met la carte en opposition & \textbf{oui} & Une référence d'opposition,
  l'heure \\
\f{preparer\_virement} & Prépare un virement vers un bénéficiaire enregistré ; vérifie le solde ;
  ne l'exécute qu'après accord & \textbf{oui} & Un récapitulatif, puis une référence \\
\f{transmettre\_conseiller} & Crée une demande de rappel avec un résumé de la conversation & oui &
  Un numéro de ticket \\
\bottomrule
\end{tabularx}
\end{center}

Le résumé transmis au conseiller est construit à partir des faits du parcours (intention, infos
recueillies, propositions) : le conseiller n'a pas à reposer les questions. C'est le
\og brouillon \fg{} de Timso AI transposé : l'humain reprend la main avec le travail déjà fait.

\section{Étage 5 --- Formuler}

\paragraph{Le gabarit d'abord.} Pour la plupart des intentions de service (carte, frais, virement,
compte), une réponse à trous, écrite à l'avance et relue, est meilleure qu'une génération : exacte,
instantanée, gratuite. Le petit LLM n'intervient que lorsque la réponse dépend de plusieurs faits à
articuler (un conseil, une explication de refus).

\paragraph{Un seul appel, à la fin.} C'est le seul endroit où un LLM intervient dans le cas
normal : une fois par tour, après que tout a été compris, décidé et calculé. L'escalade de
l'étage 1 est l'exception, comptée et affichée.

\paragraph{Ce que reçoit le petit LLM.} Une consigne système courte (registre de la banque,
vouvoiement, interdits), les \emph{faits structurés} produits par les étages 3 et 4 (propositions,
raisons, chiffres), et la question en cours. Jamais le catalogue entier, jamais l'historique brut :
le contexte reste petit, donc rapide et peu coûteux.

\paragraph{Le garde-fou.} Après génération, trois contrôles automatiques : chaque nombre du texte
doit figurer dans les faits fournis ; aucune formule interdite (\og accordé \fg, \og garanti \fg,
\og sans risque \fg) ; longueur bornée. En cas d'échec, la réponse gabarit est affichée à la place
et l'échec est compté. Le taux d'échec du garde-fou fait partie des résultats présentés.

\paragraph{Le modèle.} \f{ministral-3:3b} par défaut (3,0 Go en Q4\_K\_M) : un modèle européen,
conçu par Mistral AI, à l'aise en français --- ce qui rejoint la question de la souveraineté que
CGI a soulevée. Mais ce choix n'est pas un a priori : le chapitre~\ref{ch:modeles} décrit comment
il est confirmé ou remplacé par la mesure.

\section{Mesurer, en continu}

Chaque appel à \f{reply()} renvoie une trace : le niveau qui a compris la demande (mots-clés,
classifieur ou escalade), puis, pour chaque étage traversé, son nom, le modèle utilisé, la latence, les tokens consommés (Ollama renvoie \f{prompt\_eval\_count} et
\f{eval\_count}) et, pendant les campagnes de mesure, l'énergie. L'interface affiche cette trace en
direct dans le panneau \og frugalité \fg{} : le jury voit, pendant la démo, qu'une opposition de
carte n'a appelé aucun LLM. Le protocole de mesure de l'énergie est au chapitre~\ref{ch:evaluation}.

\section{Un exemple complet : l'appartement}

Léa (C07), 31 ans, 2\,900\eur{} nets par mois, aucun crédit en cours, 25\,000\eur{} d'épargne.

\begin{center}\small
\begin{tabularx}{\textwidth}{@{}l Y l@{}}
\toprule
\textbf{Étape} & \textbf{Ce qui se passe} & \textbf{Modèle} \\
\midrule
\multicolumn{3}{@{}l}{\textit{Tour 1 --- \og Je veux acheter un appartement \fg}} \\
Comprendre & Famille \f{credit}, intention \f{credit\_immobilier}. Les mots-clés \og acheter \fg{}
  et \og appartement \fg{} désignent une seule famille, dans une phrase courte : le filtre conclut
  seul. Entité \f{objet=appartement} & mots-clés \\
Converser & Parcours \f{credit\_immobilier}. Revenus et crédits connus par le profil ; épargne
  connue, à confirmer comme apport. Manque : le prix du bien & --- \\
Formuler & \og C'est un beau projet. J'ai noté vos revenus de 2\,900\eur{} nets par mois. Quel
  est le prix du bien que vous visez, et quelle part de vos 25\,000\eur{} d'épargne souhaitez-vous y
  consacrer ? \fg & gabarit \\
\midrule
\multicolumn{3}{@{}l}{\textit{Tour 2 --- \og Autour de 250\,000, et je peux mettre les 25\,000 \fg}} \\
Comprendre & Lu comme la réponse attendue : prix 250\,000\eur{}, apport 25\,000\eur{} & règles \\
Conseiller & Frais de notaire estimés à 7,5\,\% dans l'ancien : 18\,750\eur{}. Besoin de
  financement : 243\,750\eur{}. À un taux indicatif de 3,5\,\% sur 25 ans, mensualité d'environ
  1\,220\eur{}, soit \textbf{42,1\,\% des revenus : au-dessus des 35\,\%}. Prix maximal finançable
  avec cet apport : environ 212\,000\eur{} & règles \\
Formuler & Explique le calcul, propose trois leviers (bien moins cher, apport plus élevé,
  co-emprunteur), signale qu'un conseiller peut étudier les aides (PTZ selon conditions), propose
  un rendez-vous & \f{ministral-3:3b} \\
Garde-fou & Tous les nombres du texte figurent dans les faits ; aucune promesse & règles \\
\bottomrule
\end{tabularx}
\end{center}

Le panneau \og pourquoi \fg{} affiche alors les règles appliquées (taux d'effort maximal de
35\,\%, durée maximale de 25 ans, estimation des frais de notaire), le taux indicatif avec sa date
et sa source, et la ligne écartée : \og prêt sur 30 ans : au-delà de la durée maximale \fg. Le
calcul est vérifiable à la main par le jury.

\section{Les parcours}

Neuf parcours couvrent les neuf situations du sujet ; les familles de service restantes de
Banking77 reçoivent une fiche de réponse courte et une orientation. Le détail de chaque parcours
est en annexe~\ref{an:parcours}.

\begin{center}\small
\begin{tabularx}{\textwidth}{@{}l Y Y@{}}
\toprule
\textbf{Parcours} & \textbf{Ce que l'assistant demande} & \textbf{Comment il conclut} \\
\midrule
Carte perdue (Inès) & Quelle carte ; perte ou vol & \f{bloquer\_carte} après confirmation ; rassure
  sur la responsabilité ; propose une nouvelle carte \\
Frais (Marc) & Quel frais, si plusieurs & \f{expliquer\_frais} sur l'historique ; plafond légal ;
  si les incidents se répètent, oriente vers l'offre clientèle fragile \\
Virement (Paul) & Bénéficiaire, montant, date & \f{preparer\_virement} pas à pas ; rappel :
  la banque ne demande jamais un code \\
Épargne (Sophie) & Pour qui, horizon, montant mensuel & Propositions classées, raisons, produits
  écartés \\
Appartement (Léa) & Prix, apport & Capacité, taux d'effort, leviers, rendez-vous \\
Crédit auto (Karim) & Montant, durée, apport & Taux d'effort avec le crédit en cours, mention
  légale, rendez-vous \\
Prêt étudiant (Lucas) & Études, montant, caution & Orientation vers le prêt étudiant ; Livret Jeune
  pour épargner \\
Jargon & Le terme & Fiche du glossaire \\
Insatisfaction (Julie) & Rien : on écoute & Reconnaît la frustration, s'excuse,
  \f{transmettre\_conseiller} avec le résumé des trois contacts \\
\bottomrule
\end{tabularx}
\end{center}

\section{Conformité}

\begin{itemize}
  \item \textbf{Données personnelles.} Toutes les données sont simulées ; tout est traité en local,
        aucune donnée ne sort de la machine (E8).
  \item \textbf{Règlement européen sur l'IA.} Les systèmes destinés à évaluer la solvabilité des
        personnes physiques sont classés à haut risque. Notre assistant ne note pas la solvabilité :
        il applique des règles publiques, donne une estimation indicative, et renvoie la décision à
        un conseiller. Le modèle d'appétence n'ordonne que des produits d'épargne. Le règlement
        impose aussi d'informer l'utilisateur qu'il échange avec une IA : un bandeau permanent le
        dit.
  \item \textbf{Crédit.} La mention légale accompagne toute réponse sur un crédit.
  \item \textbf{Protection.} Toute réponse qui touche à la carte ou au virement rappelle qu'aucun
        code ne doit être communiqué --- une protection simple contre l'arnaque au faux conseiller,
        qui vise en particulier les clients âgés.
\end{itemize}

\chapter{Les modèles, et comment on les choisit}
\label{ch:modeles}

\section{Les modèles retenus}

\begin{center}\small
\begin{tabularx}{\textwidth}{@{}L{2.3cm} L{4.3cm} l l Y@{}}
\toprule
\textbf{Rôle} & \textbf{Modèle} & \textbf{Paramètres} & \textbf{Taille} & \textbf{Où il tourne} \\
\midrule
Filtre d'entrée (étage 1) & Lexique de mots-clés par famille & --- & quelques Ko & CPU \\
Encodeur (étage 1) & \f{multilingual-e5-small} (intfloat) & 118 M & $\approx$\,0,47 Go & CPU ou GPU \\
Tête de classification & Régression logistique & $\approx$\,38\,000 & $<$ 1 Mo & CPU \\
Base de comparaison & TF-IDF + régression logistique & --- & quelques Mo & CPU \\
Formulation (étage 5) & \f{ministral-3:3b}, Q4\_K\_M & $\approx$\,3 G & 3,0 Go & GPU (RTX 3060) \\
Escalade & \f{ministral-3:14b}, Q4\_K\_M & $\approx$\,14 G & 9,1 Go & GPU, à la demande \\
Appétence (étage 3) & Régression logistique & $\approx$\,50 & $<$ 1 Mo & CPU \\
Préparation, hors ligne & \f{ministral-3:14b} ou plus gros, via vLLM & --- & --- & DeepEirb, une fois \\
\bottomrule
\end{tabularx}
\end{center}

La tête de classification compte environ 38\,000 poids : 99 classes (14 familles et 85
intentions) fois 385 (les 384 dimensions du vecteur, plus un biais). Les tailles des modèles
Ollama sont celles de la bibliothèque Ollama ; celles des modèles Ministral 3 incluent un encodeur
d'images que nous n'utilisons pas.

\begin{attention}
\textbf{La mémoire de la RTX 3060 est de 12 Go.} \f{ministral-3:3b} (3,0 Go) et
\f{ministral-3:14b} (9,1 Go) ne tiennent pas ensemble avec leurs caches de contexte. Trois options,
par ordre de préférence :
\begin{enumerate}
  \item l'escalade est envoyée au serveur vLLM de DeepEirb, qui sert le gros modèle à toutes les
        équipes ; le poste ne garde que le petit ;
  \item l'escalade utilise \f{ministral-3:8b} (6,0 Go), qui tient à côté du 3B ;
  \item un seul modèle est chargé à la fois ; l'escalade, rare, accepte quelques secondes de
        rechargement.
\end{enumerate}
Dans tous les cas, le contexte est réglé au plus juste (\f{num\_ctx} à 4\,096 au lieu de 16\,384) :
moins de mémoire et des réponses plus rapides, comme l'a rappelé Denis Barthou.
\end{attention}

\section{\og Le plus petit modèle qui suffit \fg, rendu opérationnel}

La règle commune ne dit pas comment on sait qu'un modèle \og suffit \fg. Nous le décidons avant
de mesurer, pour ne pas choisir après coup le résultat qui nous arrange. La démarche est celle
qu'un organisateur nous a conseillée : partir d'un modèle qui marche, puis descendre.

\begin{cle}
\textbf{Protocole.} On commence par un modèle qui fait clairement l'affaire : il fixe la référence
de qualité. On descend ensuite en taille, candidat par candidat, sur le même jeu et avec le même
seuil fixé à l'avance. \textbf{On retient le plus petit qui tient encore le seuil.} Un modèle plus
gros n'est retenu que si tous les plus petits échouent --- c'est la définition exacte de \og un
modèle plus gros seulement si le gain est mesuré \fg. Toute l'échelle est évaluée : la courbe
qualité, taille et énergie qui en sort va dans le support.
\end{cle}

\paragraph{Pour l'encodeur.} Candidats : \f{paraphrase-multilingual-MiniLM-L12-v2} et
\f{multilingual-e5-small} (118 M chacun), puis \f{multilingual-e5-base} (278 M). Seuil : rester à
moins d'un point de F1 macro du meilleur candidat testé, au niveau des familles, sur le test
français.

\paragraph{Pour la formulation.} Référence : \f{ministral-3:14b} (9,1 Go, déjà installé), ou un
modèle d'environ 30B sur DeepEirb si Denis Barthou en sert un. Puis, en descendant :
\f{gemma4:e2b} (7,2 Go, déjà installé), \f{ministral-3:8b} (6,0 Go), \f{ministral-3:3b}
(3,0 Go), \f{qwen3:1.7b} (1,4 Go, réflexion désactivée), \f{gemma3:1b} (0,8 Go),
\f{qwen3:0.6b} (0,5 Go) et \f{gemma3:270m} (0,3 Go). Jeu : 30 situations sous forme de faits structurés, tirées des parcours. Seuil :
\textbf{zéro nombre inventé après garde-fou} et au moins 27 réponses sur 30 jugées correctes et
claires en français par deux membres de l'équipe. Notre pari est que \f{ministral-3:3b} sera le plus
petit à tenir le seuil ; la mesure peut descendre plus bas.

\paragraph{Les réglages de génération.} Température basse (0,2), au plus 200 tokens produits,
réflexion désactivée sur les modèles qui en ont une, modèle gardé en mémoire pendant la démo
(\f{keep\_alive}) pour ne pas payer le chargement devant le jury.

\section{Le gros modèle, à l'instant précis}

CGI a encouragé l'usage d'un gros modèle \og à un instant T, pour un besoin spécifique \fg. Nous en
faisons deux usages, et seulement deux.

\paragraph{Hors ligne, une fois : préparer les données.} Traduire les 13\,083 phrases de Banking77
et générer les exemples des nouvelles intentions est un travail de qualité, fait une seule fois.
C'est le bon endroit pour un gros modèle : son coût est payé une fois, mesuré avec CodeCarbon, et
\emph{amorti} sur toutes les requêtes futures. Le pitch donne les deux chiffres : l'énergie totale
de la préparation, et l'énergie par requête en service.

\paragraph{En ligne, rarement : lever un doute.} Quand le classifieur reste sous le seuil après une
question de clarification. Le taux d'escalade est mesuré sur les scénarios et affiché ; une
escalade fréquente signalerait un classifieur insuffisant, pas un besoin de gros modèle.

\section{La traduction de Banking77}

\paragraph{La méthode.} Traduction par lots avec une consigne et un glossaire : adapter au contexte
bancaire français plutôt que traduire mot à mot.

\begin{center}\small
\begin{tabular}{@{}l l@{}}
\toprule
\textbf{Anglais (Banking77)} & \textbf{Français attendu} \\
\midrule
\emph{top up} & alimenter son compte, recharger \\
\emph{transfer} & virement \\
\emph{direct debit} & prélèvement \\
\emph{statement} & relevé de compte \\
\emph{PIN} & code confidentiel \\
\emph{ATM} & distributeur \\
\emph{card payment declined} & paiement par carte refusé \\
\emph{exchange rate} & taux de change \\
\bottomrule
\end{tabular}
\end{center}

\paragraph{Le contrôle.} Contrôles automatiques sur tout le jeu (la sortie est bien en français,
le rapport de longueur avec l'original est plausible, pas de ligne vide) ; relecture humaine de
100 phrases tirées au hasard, avec le taux de phrases acceptables reporté dans la fiche de données.

\paragraph{Le repli.} Si DeepEirb n'est pas accessible lundi soir, la traduction passe par
\f{Helsinki-NLP/opus-mt-en-fr}, un modèle de traduction dédié d'environ 75 millions de paramètres,
très rapide sur la RTX 3060. Il connaît moins bien le jargon bancaire : le glossaire est alors
appliqué en post-traitement, et la relecture porte sur le jeu de test en priorité.

\section{Ce qu'il faut demander lundi soir}

\begin{itemize}
  \item \textbf{À Denis Barthou} : l'accès SSH à DeepEirb, et un modèle servi par vLLM pour la
        préparation et l'escalade (\f{ministral-3:14b} au minimum, un modèle plus gros si
        disponible pour la traduction).
  \item \textbf{Sur les postes, en filaire} : télécharger \f{ministral-3:3b}, \f{qwen3:1.7b},
        \f{gemma3:1b} et \f{multilingual-e5-small}. Les slides demandent d'éviter les gros
        téléchargements en wifi.
  \item \textbf{Aux organisateurs} : la version enrichie de Bank Marketing, et la confirmation de
        la machine de démonstration (DeepEirb ou poste de la salle).
\end{itemize}

\chapter{Évaluer et mesurer}
\label{ch:evaluation}

La qualité technique vaut 20\,\% et s'intitule \og résultats mesurés et bien évalués \fg ; la
frugalité vaut 20\,\% et se mesure. Ce chapitre fixe \emph{avant} de coder ce qu'on mesure, sur
quoi, et comment. Les tableaux sont des gabarits : ils seront remplis mardi soir et présentés tels
quels.

\section{Ce qu'on mesure}

\begin{center}\small
\begin{tabularx}{\textwidth}{@{}L{3.6cm} Y L{3.2cm} l@{}}
\toprule
\textbf{Mesure} & \textbf{Comment} & \textbf{Sur quoi} & \textbf{Critère} \\
\midrule
Précision, F1 macro & scikit-learn, par intention et par famille & Tests Banking77 et nos tests &
  Technique \\
Taux de réussite & Rejeu automatique des scénarios & 30 conversations & Technique, métier \\
Échecs du garde-fou & Compteur dans la trace & Toutes les réponses générées & Explicabilité \\
Répartition par niveau & Compteur dans la trace : mots-clés, classifieur, escalade & Tests et
  scénarios & Frugalité \\
Taux d'escalade & Compteur dans la trace & Scénarios & Frugalité \\
Part sans LLM & Compteur dans la trace & Scénarios & Frugalité \\
Latence & Horloge par étage (p50, p95) & Scénarios & Frugalité \\
Énergie & CodeCarbon, par lot (section~\ref{sec:energie}) & Lots de 200 requêtes & Frugalité \\
Taille & Go sur disque, quantisation comprise & Chaque modèle chargé & Frugalité \\
Appétence & Aire sous la courbe ROC & 20\,\% de Bank Marketing tenus à part & Technique \\
\bottomrule
\end{tabularx}
\end{center}

\section{La compréhension}

\paragraph{Trois jeux de test.} (a) Le test officiel de Banking77 en anglais, 3\,080 phrases ;
(b) le même, traduit en français ; (c) nos 160 phrases françaises des 8 nouvelles intentions,
écrites à la main. Le jeu (a) rend nos chiffres comparables à l'article d'origine ; les jeux (b) et
(c) mesurent ce qui compte pour le sujet.

\begin{center}\small
\begin{tabularx}{\textwidth}{@{}Y c c c c c@{}}
\toprule
\textbf{Configuration} & \textbf{Préc. (a)} & \textbf{Préc. (b)} & \textbf{F1 famille (b+c)} &
  \textbf{Latence p95} & \textbf{Taille} \\
\midrule
\textit{Référence : BERT affiné (article)} & \textit{93,66\,\%} & --- & --- & --- & \textit{1,4 Go} \\
Mots-clés seuls, sur la part acceptée & \amesurer & \amesurer & \amesurer & \amesurer & Ko \\
TF-IDF + régression logistique & \amesurer & \amesurer & \amesurer & \amesurer & \amesurer \\
MiniLM-L12 + régression logistique & \amesurer & \amesurer & \amesurer & \amesurer & 0,47 Go \\
\textbf{e5-small + régression logistique} & \amesurer & \amesurer & \amesurer & \amesurer &
  \textbf{0,47 Go} \\
e5-base + régression logistique & \amesurer & \amesurer & \amesurer & \amesurer & 1,1 Go \\
\f{qwen3:1.7b}, sans exemple & \amesurer & \amesurer & \amesurer & \amesurer & 1,4 Go \\
\f{ministral-3:14b}, sans exemple & \amesurer & \amesurer & \amesurer & \amesurer & 9,1 Go \\
\bottomrule
\end{tabularx}
\end{center}

Pour les mots-clés, on donne aussi la part des messages qu'ils acceptent : l'organisateur l'estime
à environ 30\,\%, nous la mesurons. La taille de la référence correspond à BERT-Large (340 millions de paramètres en float32). Les
deux dernières lignes classent avec un LLM, la liste des intentions dans la consigne : elles
montrent ce que coûterait la compréhension \og tout LLM \fg. Le message attendu : un encodeur de
118 millions de paramètres fait aussi bien, ou presque, pour une fraction du coût.

\paragraph{Le seuil $\tau$.} Sur un jeu de validation (10\,\% de l'entraînement, tenus à part), on
trace la précision des messages acceptés en fonction de la part de messages acceptés, pour $\tau$
de 0,30 à 0,95. On retient le plus petit $\tau$ qui donne au moins 97\,\% de précision. La courbe
elle-même va dans le support : elle montre qu'on a \emph{choisi} le compromis.

\section{Les scénarios de bout en bout}

Un scénario est une conversation complète, avec un client donné et des attentes vérifiables
automatiquement.

\begin{code}
\footnotesize
\begin{verbatim}
# eval/scenarios/sophie_epargne.yaml
client: C03
tours:
  - client: "je voudrais mettre de cote pour mes enfants"
    attendu: {famille: epargne, demande: [horizon]}
    interdit: {demande: [age, enfants]}        # deja connus par le profil
  - client: "dans 10 ans pour leurs etudes, 100 euros par mois"
    attendu: {propose: [livret_a], ecarte: [lep]}
    interdit: {nombre_invente: true}
\end{verbatim}
\end{code}

\paragraph{Le jeu.} Trente conversations : les neuf parcours, chacun en trois variantes
(formulation différente, autre client, réponse en plusieurs tours), plus trois demandes hors
périmètre. Chaque membre écrit les scénarios de son module ; Sohayb maintient le programme qui les
rejoue. Le jeu est rejoué à chaque intégration : \textbf{un scénario qui passait et qui échoue
bloque la fusion}.

\paragraph{Ce qui est vérifié.} La famille reconnue ; les questions posées, et surtout celles qui
\emph{ne doivent pas} l'être (une information déjà connue) ; les produits proposés et écartés ;
l'action déclenchée et sa confirmation ; le passage de main ; l'absence de nombre inventé.

\section{La formulation}

Deux indicateurs. \textbf{Le taux d'échec du garde-fou}, automatique, sur toutes les réponses
générées. \textbf{Une évaluation humaine courte} : pour chaque modèle candidat, du 14B au plus petit, 30 réponses notées
par deux membres sur trois critères (exact, clair, ton juste), en oui ou non, \textbf{champs
facultatifs} --- la leçon du formulaire de Timso AI.

\section{La frugalité : trois configurations}

Le cœur de la démonstration de frugalité est une comparaison sur le même jeu de scénarios.

\begin{center}\small
\begin{tabularx}{\textwidth}{@{}l Y@{}}
\toprule
\textbf{Configuration} & \textbf{Ce qu'elle fait} \\
\midrule
A --- Tout LLM & \f{ministral-3:14b} reçoit le message, le profil et le catalogue dans sa
  consigne, et fait tout : comprendre, décider, rédiger. C'est ce que ferait une équipe pressée. \\
B --- Cascade & Notre architecture complète. \\
C --- Cascade sans LLM & La même, avec des gabarits à la place du petit LLM : la borne basse en
  énergie. \\
\bottomrule
\end{tabularx}
\end{center}

\begin{center}\small
\begin{tabular}{@{}l c c c c c c@{}}
\toprule
& \textbf{Réussite} & \textbf{Nombres inventés} & \textbf{Wh / 1\,000 req.} &
  \textbf{Latence p95} & \textbf{Go chargés} & \textbf{Sans LLM} \\
\midrule
A & \amesurer & \amesurer & \amesurer & \amesurer & 9,1 & 0\,\% \\
B & \amesurer & \amesurer & \amesurer & \amesurer & 3,5 & \amesurer \\
C & \amesurer & \amesurer & \amesurer & \amesurer & 0,5 & 100\,\% \\
\bottomrule
\end{tabular}
\end{center}

Ce qu'on s'attend à montrer, et qu'il faudra confirmer : B atteint la qualité de A, voire mieux
parce que ses décisions sont des règles ; B consomme une fraction de l'énergie de A ; C coûte
encore moins mais rédige moins naturellement --- ce qui justifie, mesure à l'appui, la présence
d'un petit LLM. Si la mesure dit autre chose, c'est la mesure qu'on présente.

\section{Le protocole de mesure de l'énergie}
\label{sec:energie}

CodeCarbon mesure l'énergie de la machine entière (processeur, mémoire, GPU Nvidia par NVML). Une
requête dure trop peu pour être mesurée seule : on mesure par lots.

\begin{enumerate}
  \item Charger les modèles et faire une requête de chauffe (le chargement n'est pas compté).
  \item Mesurer 60 secondes à vide : puissance de repos $P_0$.
  \item Rejouer un lot de $N = 200$ requêtes tirées des scénarios ; énergie totale $E$, durée $t$.
  \item Énergie par requête : $(E - P_0 \times t) / N$.
  \item Répéter trois fois ; présenter la moyenne et l'écart.
\end{enumerate}

\begin{code}
\footnotesize
\begin{verbatim}
from codecarbon import OfflineEmissionsTracker
suivi = OfflineEmissionsTracker(country_iso_code="FRA", measure_power_secs=1,
                                project_name="cascade-B", log_level="error")
suivi.start()
for message, client in lot:
    reply(message, client, {})
suivi.stop()
energie_kwh = suivi.final_emissions_data.energy_consumed
\end{verbatim}
\end{code}

La version \og hors ligne \fg{} de CodeCarbon n'interroge aucun service extérieur : cohérent avec
la règle. On présente d'abord l'énergie en wattheures, puis l'équivalent CO\textsubscript{2}
calculé par CodeCarbon pour le réseau français. Sur Jetson, si le temps le permet, la puissance est
lue avec \f{tegrastats}.

\chapter{L'interface et la démonstration}
\label{ch:demo}

\section{L'interface}

Une seule page Streamlit, trois zones. La conversation à gauche ; à droite, le profil du client, le
panneau \og pourquoi \fg{} et le panneau \og frugalité \fg. Le jury doit pouvoir lire la décision et
son coût sans qu'on ait à les commenter.

\begin{figure}[H]
\centering
\begin{adjustbox}{max width=\textwidth}
\begin{tikzpicture}[x=1cm,y=1cm,font=\scriptsize]
  % Cadre
  \draw[darkgray,rounded corners=2pt,fill=white] (0,0) rectangle (15.4,-9.4);
  \fill[mainblue] (0,0) rectangle (15.4,-0.7);
  \node[white,anchor=west,font=\small\bfseries] at (0.25,-0.35) {Assistant bancaire --- prototype};
  \node[white,anchor=east] at (15.2,-0.35)
    {[ Client : Léa (C07) $\blacktriangledown$ ] \quad [ Mots simples ] \quad [ A+ ]};
  \fill[ext!70] (0,-0.7) rectangle (15.4,-1.15);
  \node[anchor=west] at (0.25,-0.92) {Vous échangez avec un assistant automatique. Il oriente ;
    un conseiller décide.};
  % Conversation
  \draw[darkgray!60] (0.25,-1.4) rectangle (9.2,-9.15);
  \node[anchor=north west,font=\scriptsize\bfseries,text=darkgray] at (0.35,-1.45) {Conversation};
  \node[acteur,anchor=east,text width=4.6cm,align=left] at (9.0,-2.25)
    {Je veux acheter un appartement};
  \node[etage,anchor=west,text width=6.2cm,align=left] at (0.45,-3.35)
    {C'est un beau projet. J'ai noté vos revenus de 2\,900\eur{} nets par mois. Quel est le prix
    du bien que vous visez, et quelle part de vos 25\,000\eur{} d'épargne souhaitez-vous y
    consacrer ?};
  \node[acteur,anchor=east,text width=4.6cm,align=left] at (9.0,-4.55)
    {Autour de 250\,000, et je peux mettre les 25\,000};
  \node[etage,anchor=west,text width=6.2cm,align=left] at (0.45,-5.95)
    {Avec ce prix, la mensualité serait d'environ 1\,220\eur{}, soit 42\,\% de vos revenus : au-delà
    des 35\,\% retenus pour un crédit immobilier. Avec votre apport, un bien autour de
    212\,000\eur{} resterait dans cette limite. \underline{Taux d'effort} \dots};
  \node[module,anchor=west,text width=6.2cm,align=left] at (0.45,-7.45)
    {\textbf{Action proposée} : prendre rendez-vous avec un conseiller \quad [ Oui ] [ Non ]};
  \draw[darkgray!60] (0.45,-8.45) rectangle (9.0,-8.95);
  \node[anchor=west,text=darkgray] at (0.55,-8.7) {Écrivez votre message\dots};
  % Profil
  \draw[darkgray!60] (9.45,-1.4) rectangle (15.15,-3.45);
  \node[anchor=north west,font=\scriptsize\bfseries,text=darkgray] at (9.55,-1.45) {Profil};
  \node[anchor=north west,text width=5.4cm,align=left] at (9.55,-1.8)
    {Léa, 31 ans, cadre\\
     Aucun crédit en cours\\
     Revenus : \textit{2\,900\eur{} nets par mois}\\
     Épargne : \textit{25\,000\eur{}}\quad{\tiny\textit{(italique : simulé)}}};
  % Pourquoi
  \draw[darkgray!60] (9.45,-3.6) rectangle (15.15,-6.95);
  \node[anchor=north west,font=\scriptsize\bfseries,text=darkgray] at (9.55,-3.65) {Pourquoi};
  \node[anchor=north west,text width=5.4cm,align=left] at (9.55,-4.0)
    {Intention : crédit immobilier\\
     Mots-clés : acheter, appartement\\
     Taux d'effort $\leq$ 35\,\% : \textcolor{kored}{\textbf{42,1\,\%}}\\
     Durée $\leq$ 25 ans\\
     Taux indicatif : 3,5\,\% (source, date)\\
     Proposé : rendez-vous conseiller\\
     Écarté : prêt sur 30 ans};
  % Frugalité
  \draw[darkgray!60] (9.45,-7.1) rectangle (15.15,-9.15);
  \node[anchor=north west,font=\scriptsize\bfseries,text=darkgray] at (9.55,-7.15) {Frugalité};
  \node[anchor=north west,text width=5.4cm,align=left] at (9.55,-7.5)
    {Comprendre : mots-clés, \dots{} ms\\
     Conseiller : règles, \dots{} ms\\
     Formuler : \texttt{ministral-3:3b}, \dots{} s\\
     Total : \dots{} s ; \dots{} Wh ; 3,5 Go};
\end{tikzpicture}
\end{adjustbox}
\caption{Maquette de l'écran. Les valeurs sont des exemples ; les points de suspension sont
remplis en direct par les mesures de la trace.}
\label{fig:maquette}
\end{figure}

\paragraph{Accessibilité.} Taille de texte réglable, contraste élevé, mode \og mots simples \fg,
termes du glossaire explicables d'un clic. C'est modeste, mais c'est ce que La Banque Postale a
demandé de ne jamais oublier, et ça se montre en dix secondes.

\paragraph{Les messages de démo.} Une liste latérale contient les messages exacts du script : un
clic les envoie. Pas de faute de frappe devant le jury, pas de temps perdu à taper.

\section{Le script de la démonstration}

Six minutes en direct, sur la machine fournie. Une personne pilote, une autre commente.

\begin{center}\small
\begin{tabularx}{\textwidth}{@{}l L{3.1cm} Y L{2.4cm}@{}}
\toprule
\textbf{Temps} & \textbf{Scène} & \textbf{Ce que le jury voit} & \textbf{Critère} \\
\midrule
0:00 & Inès : \og J'ai perdu ma carte hier soir \fg & L'assistant demande confirmation, met la carte
  en opposition, donne une référence et rassure. Panneau frugalité : \textbf{compris par les
  mots-clés, aucun LLM appelé}. &
  Métier, frugalité \\
1:00 & Léa : l'appartement & Il ne redemande pas les revenus ; il calcule, refuse avec le calcul,
  propose des leviers. Panneau pourquoi : la règle des 35\,\%, la ligne écartée. & Explicabilité \\
2:30 & Sophie puis Lucas : \og Je voudrais mettre de côté \fg & La même phrase, deux clients, deux
  conseils différents, chacun avec ses raisons. & Personnalisation \\
3:30 & Marc : les 8\eur{} de frais & Il retrouve les frais dans l'historique, les explique, cite
  le plafond légal, oriente vers l'offre clientèle fragile. & Métier, inclusion \\
4:30 & Julie : \og Votre appli ne marche jamais ! \fg & Il reconnaît la frustration, ne se
  justifie pas, transmet au conseiller un résumé des trois contacts. & Métier \\
5:15 & Le bilan & Panneau frugalité cumulé ; graphique des configurations A, B et C. & Frugalité \\
\bottomrule
\end{tabularx}
\end{center}

Paul et le virement en mots simples sont dans la vidéo de secours et servent de réponse si le jury
demande \og et pour un client âgé ? \fg.

\section{Le plan de secours}

\begin{itemize}
  \item \textbf{Une vidéo} de tout le script, enregistrée mercredi matin sur la machine de démo.
  \item \textbf{Un mode gabarits} : un interrupteur coupe le petit LLM. La démo continue, moins
        fluide mais exacte. C'est aussi la configuration C, donc ce n'est pas un aveu d'échec.
  \item \textbf{Une vérification réelle} d'Ollama au lancement de l'interface (pas seulement la
        présence d'une configuration) et des modèles préchargés avant l'entrée du jury.
  \item \textbf{Une répétition sur la machine de démo} mercredi entre 12h et 13h15.
\end{itemize}

\chapter{L'organisation de l'équipe}
\label{ch:organisation}

\section{Les principes}

\begin{itemize}
  \item \textbf{Un module, un responsable, un suppléant.} Le responsable décide dans son module ;
        le suppléant connaît le code assez pour le reprendre si besoin, et relit chaque fusion.
  \item \textbf{Des contrats figés lundi soir.} Les modules communiquent par des structures de
        données définies ensemble, une fois. Chacun développe contre des \emph{bouchons} qui
        respectent le contrat : personne n'attend personne.
  \item \textbf{\f{main} tourne toujours.} À tout moment, \f{streamlit run app.py} doit lancer la
        démo depuis \f{main}.
  \item \textbf{Deux intégrations mardi}, à 12h30 et à 18h, avec le jeu de scénarios comme juge.
  \item \textbf{D'abord une solution sûre, ensuite les idées.} Jusqu'à J2, on construit ce qui est
        prévu. Après, on essaie les idées nouvelles, même exotiques : c'est ce que l'organisateur
        nous a encouragés à faire, une fois la base solide. Et on se fait plaisir.
\end{itemize}

\section{Les trois modules}

Le découpage suit la cascade, pas les filières. Chaque module contient une part de données, une
part de modèle et une part de logiciel : chacun touche à l'IA, au code et à la mesure, et chacun
peut présenter son module au jury avec des chiffres.

\begin{center}\small
\begin{tabularx}{\textwidth}{@{}L{3.3cm} Y Y Y@{}}
\toprule
& \textbf{Module 1 --- Comprendre et mesurer} & \textbf{Module 2 --- Converser et formuler} &
  \textbf{Module 3 --- Conseiller et agir} \\
\midrule
\textbf{Responsable} & Younes & Yahya & Sohayb \\
\textbf{Suppléant} & Yahya & Sohayb & Younes \\
\textbf{Étages} & 1, et la mesure transversale & 2, 5, l'escalade, l'interface & 3 et 4 \\
\textbf{Part données} & Banking77 en français, exemples des 8 intentions & Glossaire, gabarits &
  Clients simulés, catalogue \\
\textbf{Part modèle} & Mots-clés, encodeur, têtes, seuil $\tau$ & Choix du petit LLM, garde-fou & Modèle
  d'appétence \\
\textbf{Part logiciel} & \f{comprendre()}, instrumentation, bancs de mesure & \f{reply()},
  automate de dialogue, Streamlit & Moteur de règles, outils, rejeu des scénarios \\
\textbf{Dans le pitch} & Les résultats et la frugalité & Le problème, l'architecture, la démo &
  L'explicabilité et la conformité \\
\bottomrule
\end{tabularx}
\end{center}

\subsection{Module 1 --- Comprendre et mesurer (Younes)}

Le module où se trouve le calcul lourd --- inférence par lots sur DeepEirb, bancs de mesure
d'énergie --- et aussi un vrai travail de modélisation.

\begin{itemize}
  \item Obtenir l'accès à DeepEirb ; traduire Banking77 (vLLM ou, à défaut, Marian sur la 3060) ;
        générer et faire relire les exemples des 8 nouvelles intentions ; écrire la fiche de
        données.
  \item Construire le filtre par mots-clés et sa règle de confiance ; entraîner l'encodeur et les
        deux têtes (famille, intention) ; choisir $\tau$ par balayage ;
        écrire les règles d'extraction des entités ; exposer \f{comprendre(message, etat)}.
  \item Remplir le tableau de compréhension du chapitre~\ref{ch:evaluation}.
  \item Écrire l'instrumentation (\f{src/mesure/frugalite.py}) et le banc A/B/C ; produire le
        graphique du pitch.
  \item En bonus : faire tourner l'étage 1 et les règles sur la Jetson.
\end{itemize}
\textbf{Fichiers} : \f{data/processed/banking77\_fr/}, \f{data/processed/intentions\_fr.csv},
\f{src/comprendre/}, \f{src/mesure/}, \f{eval/bench\_nlu.py}, \f{eval/bench\_frugalite.py}.

\subsection{Module 2 --- Converser et formuler (Yahya)}

Le module de l'orchestration : il appelle les deux autres et porte l'interface.

\begin{itemize}
  \item Écrire l'automate de dialogue et les neuf parcours YAML ; écrire \f{reply()}, le point
        d'entrée unique de la cascade.
  \item Écrire les gabarits ; brancher le petit LLM via Ollama ; écrire le garde-fou ; dérouler le
        protocole de choix du modèle de formulation.
  \item Brancher l'escalade (DeepEirb ou Ollama) ; écrire le glossaire et sa recherche en mémoire.
  \item Construire l'interface : conversation, profil, panneaux \og pourquoi \fg{} et \og
        frugalité \fg, mode mots simples, messages de démo ; enregistrer la vidéo de secours.
\end{itemize}
\textbf{Fichiers} : \f{src/assistant.py}, \f{src/converser/}, \f{app.py},
\f{data/processed/glossaire.json}.

\subsection{Module 3 --- Conseiller et agir (Sohayb)}

Le module des décisions : il doit être juste, testé et lisible, ce qui en fait un module de génie
logiciel autant que d'IA.

\begin{itemize}
  \item Écrire l'adaptateur de la base clients et l'enrichissement simulé ; créer les huit clients
        de démonstration.
  \item Écrire le catalogue JSON avec sources et dates ; écrire le moteur de règles (éligibilité,
        pertinence, classement, raisons, produits écartés) et les calculs de crédit.
  \item Entraîner et évaluer le modèle d'appétence ; exposer ses trois facteurs principaux.
  \item Écrire les cinq outils, la confirmation et le journal.
  \item Écrire le programme de rejeu des scénarios ; le lancer à chaque intégration ; tenir le dépôt
        Git propre jusqu'au dépôt final.
\end{itemize}
\textbf{Fichiers} : \f{src/conseiller/}, \f{data/processed/clients.csv},
\f{data/processed/catalogue.json}, \f{eval/run\_scenarios.py}, \f{tests/}.

\subsection{Ce qui est commun}

Chacun écrit les scénarios de son module (au moins dix). Chacun écrit les diapositives de son
module. Les contrats sont écrits ensemble, lundi soir, et ne changent plus qu'à l'unanimité.

\section{Les contrats entre modules}

Un seul fichier, \f{src/contrats.py}, importé par tous. Les noms de champs sont en français et sans
accents.

\begin{code}
\footnotesize
\begin{verbatim}
from dataclasses import dataclass, field

@dataclass
class Comprehension:                  # rendu par le module 1
    famille: str                      # "credit", "carte_perdue", ...
    intention: str                    # "credit_immobilier", "lost_or_stolen_card", ...
    confiance: float                  # probabilite de la famille, entre 0 et 1
    alternatives: list                # [("virement", 0.21), ...] pour la clarification
    entites: dict                     # {"montant": 250000, "lien": "fils", "objet": "voiture"}
    insatisfaction: bool = False
    modele: str = "e5-small+lr"       # ou "mots-cles", ou "ministral-3:14b" (escalade)
    exemple_proche: str = ""          # phrase d'entrainement la plus proche (explication)

@dataclass
class Proposition:                    # rendu par le module 3
    produit_id: str                   # "livret_a"
    raisons: list                     # ["Lucas a 20 ans : Livret Jeune possible (12-25 ans)"]
    sources: list                     # URL publiques du catalogue
    score: float = 0.0

@dataclass
class Conseil:                        # rendu par le module 3
    propositions: list                # [Proposition], triees
    ecartes: list                     # [("lep", "revenu au-dessus du plafond")]
    manque: list                      # ["prix_bien"] : informations a demander
    alertes: list                     # ["taux d'effort 42,1 % > 35 %"]
    faits: dict                       # tous les nombres que le texte a le droit de citer

@dataclass
class Action:                         # rendu par le module 3
    outil: str                        # "bloquer_carte"
    parametres: dict
    confirmee: bool
    resultat: dict = field(default_factory=dict)   # {"reference": "OPP-..."}

@dataclass
class Etape:                          # une ligne du panneau frugalite (module 1)
    nom: str; modele: str; latence_ms: float; tokens: int = 0

@dataclass
class Reponse:                        # rendu par le module 2, affiche par l'interface
    texte: str
    comprehension: Comprehension
    conseil: Conseil | None
    action: Action | None
    escalade: bool
    trace: list                       # [Etape]

def reply(message: str, client_id: str, etat: dict) -> tuple[Reponse, dict]: ...
\end{verbatim}
\end{code}

\section{L'arborescence du dépôt}

\begin{code}
\footnotesize
\begin{verbatim}
hackathon-IA/
|-- app.py                       interface Streamlit                      (M2)
|-- src/
|   |-- contrats.py              structures partagees                     (tous)
|   |-- assistant.py             reply() : la cascade                     (M2)
|   |-- comprendre/              nlu.py, entites.py, tetes/               (M1)
|   |-- converser/               dialogue.py, parcours/*.yaml,
|   |                            formuler.py, garde_fou.py, glossaire.py  (M2)
|   |-- conseiller/              clients.py, catalogue.py, regles.py,
|   |                            appetence.py, outils.py                  (M3)
|   `-- mesure/                  frugalite.py                             (M1)
|-- data/
|   |-- raw/                     Banking77, Bank Marketing (ignore par git)
|   `-- processed/               jeux francais, clients, catalogue, glossaire
|-- eval/
|   |-- scenarios/*.yaml         30 conversations                         (tous)
|   |-- run_scenarios.py                                                  (M3)
|   |-- bench_nlu.py, bench_frugalite.py                                  (M1)
|-- tests/                       tests unitaires des regles               (M3)
`-- rapport/                     ce document
\end{verbatim}
\end{code}

Les données téléchargées sont aujourd'hui dans \f{data/banking77/} et \f{data/bank-marketing/} :
elles passent dans \f{data/raw/}, ajouté au \f{.gitignore}. Les jeux préparés, petits et produits
par nous, sont versionnés : ils font partie du livrable.

\section{Le planning}

\begin{table}[H]
\centering\scriptsize
\renewcommand{\arraystretch}{1.25}
\begin{tabularx}{\textwidth}{@{}L{1.45cm} Y Y Y Y Y@{}}
\toprule
& \textbf{Lundi\newline 18h--20h} & \textbf{Mardi\newline 8h30--12h30} & \textbf{Mardi\newline 13h30--18h}
  & \textbf{Mardi\newline 18h--20h (option)} & \textbf{Mercredi\newline 8h30--12h} \\
\midrule
\textbf{Younes}\newline (M1) & \cellcolor{ent}Accès DeepEirb ; lancer la traduction
  & \cellcolor{ent}Mots-clés ; têtes v1 (anglais et français partiel) ; entités ; \texttt{comprendre()}
  & \cellcolor{ent}Français complet ; 8 intentions ; seuil $\tau$ ; escalade ; instrumentation
  & \cellcolor{ent}Banc A/B/C sans surveillance
  & \cellcolor{ent}Chiffres finaux ; graphique ; Jetson (bonus) \\
\textbf{Yahya}\newline (M2) & \cellcolor{tab}Parcours carte ; squelette de \texttt{app.py}
  & \cellcolor{tab}Automate ; parcours carte et frais ; gabarits ; petit LLM branché
  & \cellcolor{tab}Parcours épargne, crédits, virement, insatisfaction ; garde-fou ; panneaux
  & \cellcolor{tab}Scénarios en échec
  & \cellcolor{tab}Vidéo de secours ; finitions de l'interface \\
\textbf{Sohayb}\newline (M3) & \cellcolor{ext}8 clients ; catalogue v0
  & \cellcolor{ext}\texttt{clients.py} ; règles épargne et crédit ; opposition ; frais
  & \cellcolor{ext}Appétence ; produits écartés ; virement ; conseiller ; rejeu et scénarios
  & \cellcolor{ext}Scénarios en échec
  & \cellcolor{ext}Dépôt Git final ; README \\
\textbf{Tous} & \cellcolor{mod}Lire ce document ; contrats figés ; bouchons
  & & & & \cellcolor{mod}Diapositives ; deux répétitions ; dépôt à 11h30 \\
\midrule
\textbf{Jalon} & \textcolor{kored}{\textbf{J0}} à 20h & \textcolor{kored}{\textbf{J1}} à 12h30 :
  intégration 1 & \textcolor{kored}{\textbf{J2}} à 18h : intégration 2, gel & ---
  & \textcolor{kored}{\textbf{J3}} à 11h30 : tout est remis \\
\bottomrule
\end{tabularx}
\caption{Le planning, par demi-journée.}
\label{tab:planning}
\end{table}

\begin{center}\small
\begin{tabularx}{\textwidth}{@{}l l Y@{}}
\toprule
\textbf{Jalon} & \textbf{Quand} & \textbf{Critère de passage} \\
\midrule
J0 & Lundi, 20h & \f{src/contrats.py} écrit et fusionné ; \f{streamlit run app.py} répond avec
  les bouchons ; la traduction est lancée (ou son repli décidé). \\
J1 & Mardi, 12h30 & Intégration 1 : Inès (carte) et Marc (frais) passent de bout en bout, avec
  les vrais modules. \\
J2 & Mardi, 18h & Intégration 2 : les neuf parcours passent leurs scénarios principaux ; la trace
  est complète. \textbf{Gel fonctionnel} : après J2, on corrige et on mesure, on n'ajoute plus. \\
J3 & Mercredi, 11h30 & Tableaux de résultats remplis ; vidéo enregistrée ; deux répétitions
  chronométrées faites ; code déposé (étiquette \f{v1.0}) ; support remis. \\
\bottomrule
\end{tabularx}
\end{center}

\section{Les priorités}

\begin{center}\small
\begin{tabularx}{\textwidth}{@{}L{2.6cm} Y@{}}
\toprule
\textbf{Priorité} & \textbf{Contenu} \\
\midrule
\textbf{Indispensable} (J2) & Étage 1 mesuré en anglais et en français ; parcours carte, frais,
  épargne, appartement, crédit auto ; règles avec raisons et produits écartés ; les 8 clients ;
  petit LLM et garde-fou ; panneaux pourquoi et frugalité ; banc A/B/C. \\
\textbf{Important} & Virement en mots simples ; insatisfaction et passage de main ; jargon ; prêt
  étudiant ; appétence ; escalade. \\
\textbf{Si le temps le permet} & Idées nouvelles, même exotiques, encouragées par l'organisateur une
  fois la base sûre ; déploiement sur Jetson ; réponse dans la langue du client ;
  outils exposés en MCP ; saisie vocale locale (Whisper), utile à l'accessibilité. \\
\textbf{Hors périmètre} & Affinage d'un LLM ; connexion à un vrai système bancaire ;
  authentification ; base de données serveur. \\
\bottomrule
\end{tabularx}
\end{center}

\section{Les règles de travail}

\begin{itemize}
  \item Branches courtes \f{m1/<sujet>}, \f{m2/<sujet>}, \f{m3/<sujet>}, fusionnées dans la
        demi-journée, après relecture du suppléant et passage des scénarios.
  \item Aucune donnée brute ni aucun secret dans Git ; les modèles ne sont jamais versionnés.
  \item Un point de cinq minutes toutes les trois heures : fait, bloqué, suivant.
  \item Tout résultat chiffré est noté au moment où il est obtenu, avec la commande qui l'a
        produit : le support se construit au fil de l'eau.
  \item \textbf{Définition de \og fini \fg{} pour un module} : le code est dans \f{main}, ses
        scénarios passent, ses mesures sont notées, et sa diapositive est écrite.
\end{itemize}

\chapter{Les risques et leurs parades}
\label{ch:risques}

\begin{center}\small
\begin{longtable}{@{}L{3.6cm} c L{7.2cm} c@{}}
\toprule
\textbf{Risque} & \textbf{Prob.} & \textbf{Parade} & \textbf{Qui} \\
\midrule
\endhead
DeepEirb inaccessible ou saturé lundi soir & moyenne & Traduction par Marian sur la 3060 ;
  escalade par \f{ministral-3:8b} en local. La traduction du jeu de test passe en premier. & M1 \\
\addlinespace
Traduction de qualité médiocre & moyenne & Glossaire dans la consigne ; contrôles automatiques ;
  relecture de 100 phrases ; on présente le taux mesuré plutôt que de le cacher. & M1 \\
\addlinespace
Le petit LLM rédige mal en français & moyenne & Gabarits pour toutes les intentions de service ;
  protocole de choix qui peut monter à \f{ministral-3:8b} si le gain est mesuré. & M2 \\
\addlinespace
Les mots-clés se trompent avec assurance & faible & Seules les expressions dont la précision
  mesurée dépasse 97\,\% sont gardées ; une seule famille touchée exigée ; sinon, classifieur. & M1 \\
\addlinespace
Le classifieur confond des familles & faible & Seuil $\tau$, question de clarification, escalade ;
  taux mesurés et montrés. & M1 \\
\addlinespace
La base enrichie des organisateurs arrive tard ou dans un autre format & moyenne & Enrichissement
  simulé dès lundi ; un seul adaptateur à modifier le jour où elle arrive. & M3 \\
\addlinespace
L'intégration échoue tard & moyenne & Contrats figés à J0 ; bouchons ; deux intégrations datées ;
  scénarios comme juge. & Tous \\
\addlinespace
Mémoire de la 3060 insuffisante & moyenne & Escalade déportée sur DeepEirb ou 8B ; contexte réduit
  à 4\,096 tokens ; un seul gros modèle chargé à la fois. & M2 \\
\addlinespace
CodeCarbon ne voit pas le GPU & faible & Lecture directe de la puissance par NVML
  (\f{nvidia-smi}) ; même protocole par lots. & M1 \\
\addlinespace
La démo plante devant le jury & faible & Vidéo de secours ; mode gabarits ; vérification réelle
  d'Ollama ; répétition sur la machine de démo. & M2 \\
\addlinespace
Le temps manque & élevée & Priorités figées (chapitre~\ref{ch:organisation}) ; gel fonctionnel à
  J2 ; ce qui n'est pas indispensable est coupé sans débat. & Tous \\
\addlinespace
Un chiffre réglementaire faux dans une réponse & faible & Tout chiffre vient du catalogue, avec sa
  source et sa date ; garde-fou sur les nombres ; relecture du catalogue par un deuxième membre. &
  M3 \\
\bottomrule
\end{longtable}
\end{center}

\begin{cle}
\textbf{Le risque principal est le temps, pas la technique.} Chaque brique de la cascade est
simple ; ce qui peut échouer, c'est de vouloir tout faire. Les priorités du
chapitre~\ref{ch:organisation} ne sont pas indicatives : à J2, ce qui n'est pas indispensable et
pas fini est abandonné.
\end{cle}

\chapter{Vérification : toutes les cases sont remplies}
\label{ch:couverture}

\section{Les exigences}

Chaque exigence de la section~\ref{sec:exigences} est rattachée à un composant et à une preuve
visible par le jury.

\begin{center}\small
\begin{longtable}{@{}l L{4.6cm} L{4.6cm} l@{}}
\toprule
\textbf{N\textsuperscript{o}} & \textbf{Composant} & \textbf{Preuve} & \textbf{Module} \\
\midrule
\endhead
E1 & Étage 1 : mots-clés, encodeur multilingue, têtes, entités & Précision sur le test français ; démo &
  M1 \\
E2 & Banking77 traduit, 8 intentions, glossaire, catalogue, fiche de données & Fichiers dans
  \f{data/processed/} ; taux de relecture & M1, M2, M3 \\
E3 & Parcours YAML, questions déduites du profil & Scénarios : questions attendues et interdites
  & M2 \\
E4 & Profil, règles d'éligibilité, appétence & Démo : même phrase, deux clients & M3 \\
E5 & Moteur de règles sur un catalogue de produits simples & Panneau pourquoi & M3 \\
E6 & Adaptateur et 8 clients ancrés sur Bank Marketing & Sélecteur de client & M3 \\
E7 & Catalogue avec source et date par fiche & Annexe~\ref{an:catalogue} ; panneau pourquoi &
  M3 \\
E8 & Modèles open source servis par Ollama et vLLM, CodeCarbon hors ligne & Aucun appel
  sortant & Tous \\
E9 & Protocole : descendre en taille, garder le plus petit qui tient le seuil ; escalade rare & Tableaux du
  chapitre~\ref{ch:evaluation} & M1, M2 \\
E10 & Instrumentation, banc A/B/C & Panneau frugalité ; graphique du pitch & M1 \\
E11 & Raisons, produits écartés, sources, garde-fou & Panneau pourquoi ; taux du garde-fou &
  M2, M3 \\
E12 & Étages sans LLM en millisecondes ; petit LLM limité à 200 tokens & Latence p95 mesurée & M1, M2 \\
E13 & Outils avec confirmation : opposition, virement, conseiller & Démo Inès & M3 \\
E14 & Neuf parcours & Les 30 scénarios & Tous \\
E15 & Glossaire ; prêt étudiant ; mode mots simples & Démo Lucas ; termes cliquables & M2 \\
E16 & Script de démo, vidéo, dépôt étiqueté \f{v1.0} & Jalon J3 & Tous \\
E17 & Mots simples, taille du texte, contraste ; offre clientèle fragile & Démo Marc et Paul &
  M2, M3 \\
E18 & Mots-clés et classifieur en entrée ; un seul appel LLM, à la fin & Répartition par
  niveau ; part sans LLM & M1, M2 \\
E19 & Gel fonctionnel à J2, puis exploration & Planning ; priorités & Tous \\
\bottomrule
\end{longtable}
\end{center}

\section{Les livrables du sujet}

\begin{center}\small
\begin{tabularx}{\textwidth}{@{}L{4.2cm} Y l@{}}
\toprule
\textbf{Livrable attendu} & \textbf{Ce que nous livrons} & \textbf{Module} \\
\midrule
Données préparées en français & Banking77 traduit (13\,083 phrases), 8 intentions de conseil en
  français, glossaire, catalogue ; fiche de données décrivant méthode et contrôles & M1 \\
Modèle de compréhension des intentions & Mots-clés, puis encodeur \f{multilingual-e5-small} et
  deux têtes
  (14 familles, 85 intentions), avec seuil et escalade ; résultats sur trois jeux de test & M1 \\
Chatbot relié à une base clients simulée & Interface Streamlit, \f{reply()}, 8 clients de
  démonstration, adaptateur prêt pour la base enrichie & M2, M3 \\
Moteur de conseils personnalisés & Règles d'éligibilité et de pertinence, calculs de crédit,
  appétence, raisons et produits écartés & M3 \\
Mesures de frugalité & Énergie, latence, taille, part sans LLM ; comparaison A/B/C ; coût de
  la préparation hors ligne & M1 \\
\bottomrule
\end{tabularx}
\end{center}

\section{La démarche en quatre temps}

Les slides demandent une démarche en quatre temps ; le jury \og évalue la démarche autant que le
résultat \fg. Chaque temps a sa réponse et sa diapositive.

\begin{center}\small
\begin{tabularx}{\textwidth}{@{}l Y@{}}
\toprule
\textbf{Temps} & \textbf{Notre réponse} \\
\midrule
Cadrer & Les neuf situations du sujet ; le constat que Banking77 couvre le service et pas le
  conseil ; 14 familles et 85 intentions. \\
Préparer & Traduction outillée et contrôlée ; 8 intentions créées ; tests écrits à la main ;
  enrichissement simulé déclaré comme tel ; \f{duration} exclue. \\
Concevoir & La cascade : chaque étage ne passe la main que s'il ne sait pas faire ; le gros modèle
  à l'instant précis. \\
Démontrer & Les tests officiels, les 30 scénarios, la comparaison A/B/C, la démo en direct. \\
\bottomrule
\end{tabularx}
\end{center}

\chapter{Comparaison avec la proposition reçue}
\label{ch:comparaison}

Un ami de l'équipe nous a transmis une proposition (\f{infos/propo.md}). Elle est solide et nous
en gardons l'essentiel. Ce chapitre dit ce qu'on garde, ce qu'on change, et pourquoi.

\section{Ce qu'on garde}

\begin{itemize}
  \item \textbf{Ne pas viser les 77 intentions comme des parcours} : quelques parcours
        parfaitement fonctionnels valent mieux que 77 réponses approximatives.
  \item \textbf{La chaîne} message, intention, données manquantes, questions, profil, règles,
        recommandation, explication.
  \item \textbf{Le LLM ne décide pas} : il comprend, extrait et reformule ; la recommandation vient
        de règles explicables.
  \item \textbf{Un contrat commun} autour d'une fonction centrale, défini dès le début.
  \item \textbf{Un mode sans LLM} en cas de problème, une vidéo de secours, des répétitions.
  \item \textbf{Des clients synthétiques et un catalogue} construits par nous.
\end{itemize}

\section{Ce qu'on change}

\begin{center}\small
\begin{longtable}{@{}L{2.9cm} L{4.1cm} L{4.1cm} L{3.4cm}@{}}
\toprule
\textbf{Point} & \textbf{Proposition reçue} & \textbf{Notre version} & \textbf{Pourquoi} \\
\midrule
\endhead
Usage de Banking77 & Traduire les catégories utiles, créer quelques exemples & Traduire tout et
  évaluer sur le test officiel, en anglais et en français & La qualité technique (20\,\%) demande
  des résultats \og bien évalués \fg, comparables à une référence publique \\
\addlinespace
Preuve de frugalité & Mesurer latence et consommation & Comparer trois configurations, dont une
  \og tout LLM \fg, sur le même jeu & Sans point de comparaison, un chiffre d'énergie ne prouve
  pas \og le plus petit modèle qui suffit \fg \\
\addlinespace
Gros modèle & Absent & Hors ligne pour préparer les données ; en ligne en escalade rare, comptée &
  CGI a dit que c'est \og là que vous allez innover \fg \\
\addlinespace
Profondeur & Guider & Agir, avec confirmation (opposition, virement, transmission) & L'oral
  ouvrait explicitement cette possibilité \\
\addlinespace
Explicabilité & Afficher les règles utilisées & Raisons, produits écartés, sources datées,
  facteurs d'appétence, garde-fou mesuré & 15\,\% de la note ; le \og pourquoi pas \fg{} est ce
  qui convainc \\
\addlinespace
Valeurs de La Banque Postale & Absentes & Mots simples, glossaire, offre clientèle fragile,
  passage de main & Accessibilité et inclusion ont été mises en avant par sa directrice SI \\
\addlinespace
Bank Marketing & Profils, propension \og éventuellement \fg & Appétence mesurée, \f{duration}
  exclue (fuite documentée) & Un point de rigueur que le jury peut vérifier \\
\addlinespace
Conformité & Absente & Règlement européen sur l'IA, mention légale du crédit, transparence &
  Jury bancaire \\
\addlinespace
Répartition & Données et métier / IA et backend / interface, mesures et slides & Trois modules
  qui suivent la cascade, chacun avec données, modèle et logiciel & La personne 3 cumulait
  l'interface, les mesures et les slides ; personne n'était responsable de l'évaluation \\
\bottomrule
\end{longtable}
\end{center}

\begin{cle}
\textbf{En résumé.} La proposition reçue décrit un bon MVP. La nôtre garde son squelette et ajoute
ce qui fait gagner des points sur la grille : des résultats comparables à une référence, une preuve
de frugalité par comparaison, une explicabilité qui inclut le \og pourquoi pas \fg, et les valeurs
du client.
\end{cle}

\chapter{Le pitch, et ce qu'on fait ce soir}
\label{ch:pitch}

\section{Les quinze minutes}

Trois voix, une par module : le jury entend chaque responsable parler de ce qu'il a construit.

\begin{center}\small
\begin{tabularx}{\textwidth}{@{}l L{2.6cm} Y L{2.3cm}@{}}
\toprule
\textbf{Temps} & \textbf{Partie} & \textbf{Contenu} & \textbf{Qui} \\
\midrule
0:00 -- 1:30 & Le problème & Un client écrit à sa banque ; un Français sur quatre est client de
  La Banque Postale ; l'expérience client est sa priorité n\textsuperscript{o}\,1 ; la règle :
  open source, local, frugal. & Yahya \\
1:30 -- 4:00 & La démarche & Les quatre temps ; le constat sur Banking77 ; la cascade (une
  figure) ; le gros modèle à l'instant précis. & Yahya \\
4:00 -- 10:00 & La démonstration & Le script du chapitre~\ref{ch:demo}. & Yahya pilote,
  Sohayb commente \\
10:00 -- 12:30 & Les résultats & Compréhension sur les trois jeux de test ; scénarios ; tableau
  A/B/C ; coût de la préparation hors ligne. & Younes \\
12:30 -- 14:00 & L'explicabilité & Raisons et produits écartés ; garde-fou ; règlement européen
  sur l'IA ; mention légale. & Sohayb \\
14:00 -- 15:00 & Limites et suite & Ce qui est simulé ; ce qu'il faudrait pour 17\,000 points de
  contact. & Younes \\
\bottomrule
\end{tabularx}
\end{center}

\section{Trois messages}

\begin{enumerate}
  \item \og Notre assistant comprend souvent avec de simples mots-clés, sinon avec 118 millions de
        paramètres ; il décide avec des règles que le client peut lire ; il n'appelle un LLM
        qu'une fois, à la fin, pour mettre en mots ce qui est déjà calculé. \fg
  \item \og Nous l'avons mesuré contre une version où un LLM de 14 milliards de paramètres fait
        tout. \fg{} (suivi des chiffres)
  \item \og Chaque conseil dit pourquoi --- et pourquoi pas. \fg
\end{enumerate}

\section{Les questions probables du jury}

\begin{center}\small
\begin{longtable}{@{}L{4.6cm} L{9.8cm}@{}}
\toprule
\textbf{Question} & \textbf{Réponse préparée} \\
\midrule
\endhead
Pourquoi ne pas tout confier à un LLM ? & Parce qu'on l'a mesuré : la configuration A fait tout
  avec un LLM de 14 milliards de paramètres ; voici sa qualité, son énergie et sa latence face à
  la cascade. \\
\addlinespace
Comment mesurez-vous l'énergie ? & CodeCarbon en mode hors ligne, par lots de 200 requêtes,
  puissance de repos soustraite, trois répétitions. Une requête seule est trop courte pour être
  mesurée honnêtement. \\
\addlinespace
Et si le classifieur se trompe ? & En dessous du seuil $\tau$, il ne devine pas : il pose une
  question fermée, puis escalade. $\tau$ a été choisi par balayage pour 97\,\% de précision sur
  les messages acceptés. Voici la courbe. \\
\addlinespace
Le LLM peut-il inventer un taux ? & Il ne reçoit que des faits calculés, et chaque nombre de sa
  réponse est contrôlé contre ces faits ; en cas d'écart, la réponse gabarit s'affiche. Voici le
  taux d'échec du garde-fou. \\
\addlinespace
Pourquoi Ministral ? & Parce qu'en descendant depuis le 14B, c'est le plus petit qui tient notre
  seuil de qualité en français --- voici la courbe --- et qu'il est européen. Si un plus petit
  avait suffi, le protocole l'aurait retenu. \\
\addlinespace
Des mots-clés, en 2026 ? & Oui : ils traitent une part mesurée des demandes pour un coût quasi nul,
  et ils disent quand ils ne sont pas sûrs. Le reste passe au classifieur. Moins de LLM, plus
  d'algorithmes. \\
\addlinespace
Bank Marketing, ce sont des données portugaises de 2008 ? & Oui. Elles ancrent des profils
  réalistes et donnent un signal d'appétence faible, présenté comme tel. L'adaptateur permet de
  brancher une vraie base sans toucher au reste. \\
\addlinespace
Et la conformité ? & Données simulées et traitement local. Le règlement européen classe à haut
  risque la notation de solvabilité : nous n'en faisons pas, un conseiller décide. Transparence :
  l'utilisateur sait qu'il parle à une IA. Mention légale du crédit. \\
\addlinespace
Comment passer à 17\,000 bureaux ? & Le coût par requête est mesuré, donc extrapolable. Les étages
  1 à 4 tournent sur CPU ; seul le petit LLM a besoin d'un GPU, et un serveur partagé suffit. Les
  pistes : voix locale et réponse dans la langue du client, dans la ligne de l'outil multilingue
  annoncé pour novembre. \\
\addlinespace
Pourquoi pas d'affinage ? & Deux jours, des GPU partagés, et une tête de classification entraînée
  en secondes qui suffit. C'est aussi la conclusion d'un système en production que nous avons
  construit : la recherche bat l'affinage quand le modèle de base évolue vite. \\
\addlinespace
Que se passe-t-il hors périmètre ? & L'assistant le dit, et propose un conseiller. Il ne fait pas
  semblant. \\
\bottomrule
\end{longtable}
\end{center}

\section{Ce qu'on fait ce soir}

\begin{enumerate}
  \item Lire la page de synthèse et le chapitre~\ref{ch:organisation} (vingt minutes) ; valider ou
        amender la répartition, puis ne plus y revenir.
  \item Écrire \f{src/contrats.py} ensemble (trente minutes).
  \item Rendre le dépôt GitHub privé : il est public, et il va contenir des éléments tirés de
        supports marqués \og Interne \fg.
  \item Déplacer les données téléchargées dans \f{data/raw/} et l'ajouter au \f{.gitignore}.
  \item \textbf{Younes} : écrire à Denis Barthou pour DeepEirb et le modèle vLLM ; lancer la
        traduction, ou décider du repli.
  \item \textbf{Yahya} : squelette de \f{app.py} et de \f{reply()} avec bouchons ; parcours carte.
  \item \textbf{Sohayb} : les 8 clients et le catalogue v0 ; squelette du rejeu des scénarios.
  \item En filaire, télécharger \f{ministral-3:3b}, \f{qwen3:1.7b}, \f{gemma3:1b} et
        \f{multilingual-e5-small} sur un poste de la salle.
  \item Vérifier J0 avant de partir : \f{streamlit run app.py} répond.
\end{enumerate}


\appendix
\chapter{Taxonomie des intentions}
\label{an:taxonomie}

Les 77 intentions de Banking77, rangées en neuf familles (vérifié par script : chacune une fois et
une seule), puis les huit intentions que nous créons.

\begin{center}\footnotesize
\begin{longtable}{@{}L{3.1cm} L{11.3cm}@{}}
\toprule
\textbf{Famille} & \textbf{Intentions} \\
\midrule
\endhead
\f{carte\_perdue} (4) & \f{lost\_or\_stolen\_card}, \f{compromised\_card}, \f{card\_swallowed},
  \f{lost\_or\_stolen\_phone} \\
\addlinespace
\f{carte\_usage} (10) & \f{card\_not\_working}, \f{contactless\_not\_working}, \f{pin\_blocked},
  \f{change\_pin}, \f{activate\_my\_card}, \f{card\_about\_to\_expire}, \f{card\_acceptance},
  \f{declined\_card\_payment}, \f{virtual\_card\_not\_working}, \f{pending\_card\_payment} \\
\addlinespace
\f{carte\_commande} (12) & \f{card\_arrival}, \f{card\_delivery\_estimate},
  \f{order\_physical\_card}, \f{get\_physical\_card}, \f{getting\_spare\_card},
  \f{getting\_virtual\_card}, \f{get\_disposable\_virtual\_card}, \f{disposable\_card\_limits},
  \f{card\_linking}, \f{visa\_or\_mastercard}, \f{supported\_cards\_and\_currencies},
  \f{apple\_pay\_or\_google\_pay} \\
\addlinespace
\f{frais\_contestation} (14) & \f{extra\_charge\_on\_statement}, \f{card\_payment\_fee\_charged},
  \f{transfer\_fee\_charged}, \f{cash\_withdrawal\_charge}, \f{exchange\_charge},
  \f{top\_up\_by\_card\_charge}, \f{top\_up\_by\_bank\_transfer\_charge},
  \f{transaction\_charged\_twice}, \f{card\_payment\_not\_recognised},
  \f{direct\_debit\_payment\_not\_recognised}, \f{cash\_withdrawal\_not\_recognised},
  \f{request\_refund}, \f{Refund\_not\_showing\_up}, \f{reverted\_card\_payment?} \\
\addlinespace
\f{virement} (10) & \f{transfer\_into\_account}, \f{cancel\_transfer}, \f{pending\_transfer},
  \f{transfer\_timing}, \f{failed\_transfer}, \f{declined\_transfer}, \f{beneficiary\_not\_allowed},
  \f{transfer\_not\_received\_by\_recipient}, \f{receiving\_money},
  \f{balance\_not\_updated\_after\_bank\_transfer} \\
\addlinespace
\f{compte\_identite} (8) & \f{edit\_personal\_details}, \f{verify\_my\_identity},
  \f{why\_verify\_identity}, \f{unable\_to\_verify\_identity}, \f{verify\_source\_of\_funds},
  \f{passcode\_forgotten}, \f{terminate\_account}, \f{age\_limit} \\
\addlinespace
\f{especes} (5) & \f{atm\_support}, \f{pending\_cash\_withdrawal}, \f{declined\_cash\_withdrawal},
  \f{wrong\_amount\_of\_cash\_received}, \f{balance\_not\_updated\_after\_cheque\_or\_cash\_deposit} \\
\addlinespace
\f{change\_etranger} (6) & \f{exchange\_rate}, \f{card\_payment\_wrong\_exchange\_rate},
  \f{wrong\_exchange\_rate\_for\_cash\_withdrawal}, \f{exchange\_via\_app},
  \f{fiat\_currency\_support}, \f{country\_support} \\
\addlinespace
\f{alimenter\_compte} (8) & \f{automatic\_top\_up}, \f{pending\_top\_up}, \f{top\_up\_limits},
  \f{top\_up\_reverted}, \f{topping\_up\_by\_card}, \f{top\_up\_failed}, \f{verify\_top\_up},
  \f{top\_up\_by\_cash\_or\_cheque} \\
\midrule
\f{epargne} (1) & \f{epargne\_projet} \\
\f{credit} (3) & \f{credit\_immobilier}, \f{credit\_conso}, \f{pret\_etudiant} \\
\f{virement} (+1) & \f{faire\_virement} \\
\f{aide} (1) & \f{question\_definition} \\
\f{insatisfaction} (1) & \f{insatisfaction} \\
\f{conversation} (1) & \f{hors\_sujet} \\
\bottomrule
\end{longtable}
\end{center}

Total : 14 familles, 85 intentions. Les noms \f{Refund\_not\_showing\_up} et
\f{reverted\_card\_payment?} sont repris tels qu'ils figurent dans Banking77.

\chapter{Le catalogue de produits}
\label{an:catalogue}

Produits simples, conditions publiques. \textbf{Chaque fiche du fichier JSON porte la source
publique et la date du relevé} ; les valeurs ci-dessous sont à vérifier à la source le jour de la
saisie, et les taux réglementés, qui changent, n'y figurent volontairement pas.

\begin{center}\small
\begin{longtable}{@{}L{2.5cm} L{3.1cm} L{2.6cm} L{2.2cm} L{3.7cm}@{}}
\toprule
\textbf{Produit} & \textbf{Pour qui} & \textbf{Plafond, limites} & \textbf{Argent disponible}
  & \textbf{Convient pour} \\
\midrule
\endhead
Livret A & Tout le monde, mineurs compris ; un par personne & 22\,950\eur{} & Oui & Épargne de
  précaution ; épargne d'un enfant \\
LDDS & Majeurs domiciliés fiscalement en France ; un par personne & 12\,000\eur{} & Oui &
  Compléter le Livret A \\
LEP & Revenu fiscal de référence sous un plafond annuel (à relever) & 10\,000\eur{} & Oui &
  Revenus modestes : taux plus élevé \\
Livret Jeune & 12 à 25 ans & 1\,600\eur{} & Oui & Premier livret d'un jeune \\
PEL & Tout le monde ; versements réguliers obligatoires & 61\,200\eur{} & Non : un retrait clôt le
  plan & Projet immobilier à plusieurs années \\
CEL & Tout le monde & 15\,300\eur{} & Oui & Petit projet, complément du PEL \\
\midrule
Crédit immobilier & Majeurs & Mensualités $\leq$ 35\,\% des revenus nets, assurance comprise ;
  durée $\leq$ 25 ans (27 dans certains cas) & --- & Achat d'un logement ; décision par un
  conseiller \\
Prêt à taux zéro & Primo-accédants, sous conditions de ressources et de zone & Selon conditions
  & --- & Complément d'un crédit immobilier ; éligibilité vérifiée par un conseiller \\
Prêt personnel & Majeurs & Capacité de remboursement ; mention légale obligatoire & --- & Voiture,
  travaux, équipement \\
Prêt étudiant & Étudiants majeurs & Montant et caution selon l'offre ; remboursement différé
  possible ; version garantie par l'État sous conditions & --- & Financer ses études \\
\midrule
Offre clientèle fragile & Personnes en situation de fragilité financière & Tarif mensuel plafonné ;
  commission d'intervention limitée à 4\eur{} par opération et 20\eur{} par mois & --- &
  Inclusion ; proposée par un conseiller \\
Opposition carte (service) & Tout titulaire d'une carte & --- & Immédiat & Perte, vol, fraude \\
\bottomrule
\end{longtable}
\end{center}

\paragraph{Règle générale sur les frais.} Hors offre clientèle fragile, la commission
d'intervention est plafonnée par la loi à 8\eur{} par opération et 80\eur{} par mois (article
R312-4-1 du Code monétaire et financier). C'est un plafond, pas un prix : le montant réellement
facturé est celui de la grille tarifaire de la banque, à relever sur son site.

\chapter{Le détail des parcours}
\label{an:parcours}

P : information lue dans le profil (confirmée si simulée). Q : question posée au client.
E : entité extraite du message.

\begin{center}\small
\begin{longtable}{@{}L{2.4cm} L{4.4cm} L{3.5cm} L{4.1cm}@{}}
\toprule
\textbf{Parcours} & \textbf{Informations} & \textbf{Règles, outils} & \textbf{Conclusion} \\
\midrule
\endhead
Carte perdue & Carte concernée (P ; Q s'il y en a plusieurs) ; perte ou vol (Q) ; quand (E) &
  \f{bloquer\_carte}, avec confirmation & Référence, heure ; rassurer ; proposer une nouvelle carte ;
  rappel : ne jamais communiquer son code \\
\addlinespace
Frais & Période (E ou P) ; frais visé (Q s'il y en a plusieurs) & \f{expliquer\_frais} ; plafond
  légal & Explication datée ; si les incidents se répètent, orientation vers l'offre clientèle
  fragile ; transmission si contestation \\
\addlinespace
Virement & Bénéficiaire (P ; Q) ; montant (E ; Q) ; date (Q, immédiat par défaut) &
  \f{preparer\_virement} ; solde suffisant ; bénéficiaire enregistré & Récapitulatif, confirmation,
  référence ; mode mots simples ; rappel contre l'arnaque \\
\addlinespace
Épargne & Pour qui (P : enfants ; Q) ; horizon (Q) ; montant mensuel (Q) ; épargne existante (P)
  & Éligibilité (âge, revenu, plafonds atteints) ; pertinence (horizon court : livrets ; projet
  immobilier : PEL, CEL) ; appétence & Propositions classées avec raisons ; produits écartés \\
\addlinespace
Crédit immobilier & Revenus (P) ; crédits en cours (P) ; prix (Q) ; apport (P, confirmé) & 35\,\% ;
  25 ans ; frais de notaire estimés ; capacité & Estimation, leviers, rendez-vous ; mention
  légale \\
\addlinespace
Crédit à la consommation & Objet (E) ; montant (Q) ; durée (Q) ; revenus et crédits (P) & Taux
  d'effort, crédit en cours compris ; reste à vivre (paramètre interne) & Estimation prudente ;
  mention légale ; rendez-vous \\
\addlinespace
Prêt étudiant & Statut (P) ; études et durée (Q) ; montant (Q) ; caution possible (Q) &
  Éligibilité (majeur, étudiant) & Orientation vers le prêt étudiant ; Livret Jeune ;
  rendez-vous \\
\addlinespace
Jargon & Terme (E, par le glossaire) & Recherche dans le glossaire, en mémoire & Fiche courte ;
  proposition d'un exemple \\
\addlinespace
Insatisfaction & Contacts du jour (P) & Règle des trois contacts ; ton & Excuses ; reformulation
  du problème ; \f{transmettre\_conseiller} avec résumé \\
\bottomrule
\end{longtable}
\end{center}

\chapter{Le glossaire}
\label{an:glossaire}

Fiches courtes, écrites par nous et relues ; c'est le texte que l'assistant affiche, sans LLM.

\begin{center}\small
\begin{longtable}{@{}L{3.3cm} L{11.1cm}@{}}
\toprule
\textbf{Terme} & \textbf{Explication en mots simples} \\
\midrule
\endhead
Agios & Les intérêts et frais que la banque facture quand votre compte est à découvert. \\
Apport personnel & L'argent que vous mettez vous-même dans un achat, sans l'emprunter. \\
Assurance emprunteur & Une assurance qui rembourse tout ou partie de votre crédit si vous ne
  pouvez plus le faire (décès, invalidité\dots). Elle est presque toujours demandée pour un crédit
  immobilier. \\
Commission d'intervention & Des frais facturés quand la banque doit traiter une opération qui pose
  problème, par exemple un paiement alors que votre solde est insuffisant. La loi les plafonne. \\
Découvert autorisé & La possibilité, prévue avec votre banque, d'avoir un solde négatif jusqu'à
  un certain montant. Il coûte des agios. \\
Épargne de précaution & Une réserve d'argent disponible à tout moment pour faire face à un
  imprévu ; on conseille souvent trois à six mois de dépenses. \\
Frais de notaire & Les taxes et frais payés lors de l'achat d'un logement, en plus du prix :
  environ 7 à 8\,\% dans l'ancien, 2 à 3\,\% dans le neuf. \\
IBAN & Le numéro qui identifie un compte bancaire ; il faut le donner pour recevoir un virement. \\
LDDS & Livret de développement durable et solidaire : un livret d'épargne disponible, sans impôt
  sur les intérêts, plafonné à 12\,000\eur{}, réservé aux majeurs. \\
LEP & Livret d'épargne populaire : un livret mieux rémunéré, réservé aux personnes dont les
  revenus sont sous un plafond. \\
Livret A & Le livret d'épargne le plus courant : argent disponible à tout moment, intérêts sans
  impôt, taux fixé par l'État, plafond de 22\,950\eur{}. \\
Livret Jeune & Un livret pour les 12--25 ans, plafonné à 1\,600\eur{}. \\
Opposition & Le blocage définitif d'une carte perdue ou volée, pour que personne ne puisse s'en
  servir. \\
PEL & Plan épargne logement : on verse régulièrement pendant plusieurs années, en vue d'un projet
  immobilier ; un retrait ferme le plan. \\
Prélèvement & Un paiement demandé directement à votre banque par une entreprise que vous avez
  autorisée (énergie, téléphone\dots). \\
Reste à vivre & Ce qui vous reste chaque mois une fois payés vos charges et vos crédits. \\
TAEG & Taux annuel effectif global : le coût total d'un crédit (intérêts, frais, assurance
  obligatoire) exprimé en pourcentage par an. C'est le chiffre qui permet de comparer deux offres. \\
Taux d'endettement & La part de vos revenus qui part chaque mois dans vos crédits. Pour un crédit
  immobilier, elle ne doit pas dépasser 35\,\%. \\
\bottomrule
\end{longtable}
\end{center}

\chapter{Commandes utiles}
\label{an:commandes}

\paragraph{Ollama sur un poste de la salle} (d'après les slides de Denis Barthou).

\begin{code}
\footnotesize
\begin{verbatim}
export PATH=/opt/ollama/bin:${PATH}
export OLLAMA_MODELS=/opt/ollama/models
export OLLAMA_CONTEXT_LENGTH=4096          # au plus juste (16384 par defaut sur les slides)
ollama serve &                             # un seul serveur par machine
ollama pull ministral-3:3b                 # en filaire
ollama list
\end{verbatim}
\end{code}

\paragraph{Appel à Ollama avec mesure.} La réponse contient \f{prompt\_eval\_count},
\f{eval\_count} et les durées en nanosecondes, ce qui alimente la trace.

\begin{code}
\footnotesize
\begin{verbatim}
curl -s localhost:11434/api/chat -d '{
  "model": "ministral-3:3b",
  "messages": [{"role": "system", "content": "..."},
               {"role": "user", "content": "..."}],
  "stream": false,
  "keep_alive": "30m",
  "options": {"temperature": 0.2, "num_ctx": 4096, "num_predict": 200}
}'
# Pour un modele qui raisonne (qwen3) : ajouter  "think": false
\end{verbatim}
\end{code}

\paragraph{Serveur vLLM sur DeepEirb.} vLLM expose une API compatible OpenAI ; l'adresse et le nom
du modèle sont donnés par Denis Barthou.

\begin{code}
\footnotesize
\begin{verbatim}
curl -s http://<deepeirb>:<port>/v1/chat/completions \
  -H "Content-Type: application/json" \
  -d '{"model": "<modele>", "temperature": 0,
       "messages": [{"role": "user", "content": "Traduis en francais ..."}]}'
\end{verbatim}
\end{code}

\paragraph{Lancer l'interface et les mesures.} Depuis la racine du dépôt, une fois les modules
en place :

\begin{code}
\footnotesize
\begin{verbatim}
python3 -m venv .venv && source .venv/bin/activate
pip install sentence-transformers scikit-learn streamlit pyyaml codecarbon pandas
streamlit run app.py
python eval/run_scenarios.py          # rejoue les 30 scenarios
python eval/bench_nlu.py              # tableau de comprehension
python eval/bench_frugalite.py A B C  # comparaison des configurations
\end{verbatim}
\end{code}


\end{document}
